\documentclass[aps,preprint]{revtex4}%
\usepackage{amsfonts}
\usepackage{amsmath}
\usepackage{amssymb}
\usepackage{graphicx}%
\setcounter{MaxMatrixCols}{30}
%TCIDATA{OutputFilter=latex2.dll}
%TCIDATA{Version=5.50.0.2953}
%TCIDATA{CSTFile=revtex4.cst}
%TCIDATA{Created=Tuesday, February 12, 2008 12:18:57}
%TCIDATA{LastRevised=Tuesday, February 12, 2008 13:27:33}
%TCIDATA{<META NAME="GraphicsSave" CONTENT="32">}
%TCIDATA{<META NAME="SaveForMode" CONTENT="1">}
%TCIDATA{BibliographyScheme=Manual}
%TCIDATA{<META NAME="DocumentShell" CONTENT="Articles\SW\REVTeX 4">}
%BeginMSIPreambleData
\providecommand{\U}[1]{\protect\rule{.1in}{.1in}}
%EndMSIPreambleData
\newtheorem{theorem}{Theorem}
\newtheorem{acknowledgement}[theorem]{Acknowledgement}
\newtheorem{algorithm}[theorem]{Algorithm}
\newtheorem{axiom}[theorem]{Axiom}
\newtheorem{claim}[theorem]{Claim}
\newtheorem{conclusion}[theorem]{Conclusion}
\newtheorem{condition}[theorem]{Condition}
\newtheorem{conjecture}[theorem]{Conjecture}
\newtheorem{corollary}[theorem]{Corollary}
\newtheorem{criterion}[theorem]{Criterion}
\newtheorem{definition}[theorem]{Definition}
\newtheorem{example}[theorem]{Example}
\newtheorem{exercise}[theorem]{Exercise}
\newtheorem{lemma}[theorem]{Lemma}
\newtheorem{notation}[theorem]{Notation}
\newtheorem{problem}[theorem]{Problem}
\newtheorem{proposition}[theorem]{Proposition}
\newtheorem{remark}[theorem]{Remark}
\newtheorem{solution}[theorem]{Solution}
\newtheorem{summary}[theorem]{Summary}
\newenvironment{proof}[1][Proof]{\noindent\textbf{#1.} }{\ \rule{0.5em}{0.5em}}
\begin{document}
\preprint{}
\title[TDDFT]{Time-dependent density functional theory}
\author{Brian Weiner}
\affiliation{The Pennsylvania State University, Department of Physics, University Park, PA 16802}
\author{S.B. Trickey}
\affiliation{Quantum Theory Project Department of Physics University
of Florida Gainesville, FL 32611-8435, U.S.A.}

\pacs{PACS number}

\begin{abstract}
The usual formulation of time-dependent density functional theory
(TDDFT) relies upon the Runge-Gross argument for the existence of a
1-1 map between densities and external potentials. Only in the
multi-component plasma generalization due to Li and Tong is this
TDDFT applicable to dynamical nuclei and electrons simultaneously.
Both the original and multi-component versions depend on satisfying
certain boundary conditions and the existence of temporal Taylor
expansions, requirements that have been the subject of skepticism.
Little actual application of the multi-component scheme has
occurred. This paper, the first of a series, extends our alternative
formulation of TDDFT explicitly to the non-Born-Oppenheimer case.
The approach does not depend on the 1-1 map logic as a route to
functionals of the density. Instead, the analysis proceeds by
classification of states according to densities (as in the
constrained search approach) into fiber bundles. The densities then
are used as parameters to label states (akin to coherent state
theory), specifically by finding invariances of the fibers and
mappings among them. Formal expressions for functionals of the
density that have a unique value associated to each of these classes
of states for all times are obtained in terms of those
transformations. The equation of motion for the density follows from
the time-dependent variational principle. The subsequent two papers
treat (1) approximations to the operators of the invariance groups
to construct approximate functionals systematically and (2) a
corresponding path integral formulation of TDDFT.

\end{abstract}

\startpage{1}
\maketitle
\tableofcontents


\section{Overview and Motivation}

\noindent Time-dependent density functional theory (TDDFT) is
growing in importance. Examples include Refs.
\cite{Teil92,Stener95,Saalmann96,Tong97,Serra97a,Ullrich97,AG97,Vignale97,Calvayrac98,Reinhard98,Ullrich98,Telnov98,Tong98,Campillo99,Vasiliev99,Knospe99,Hessler}
as well as citations in surveys such as Refs.
\cite{Gross96,Casida95}. To date, however, the coupled motions of
electrons and nuclei than to electron excitations and response with
static nuclei have attracted less attention.

Though there are various other formulations of TDDFT
\cite{Bartolotti81,DebGhosh82,KD86}, the conventional form is that of Runge
and Gross \cite{RG84}, hereafter RG. They proceeded in analogy with the
original ground state Hohenberg-Kohn argument \cite{HK} by reductio ad
absurdum to the result that, for a $v$-representable time-dependent density,
under certain assumptions the map between the time-dependent external
potential and that density is invertible. For those assumptions, invertibility
establishes the existence, though not the structure, of a TDDFT. Issues of
boundary conditions on the potentials and densities and implicit assumptions
for which the RG proof is valid arose shortly \cite{Xu85,Dhara87}. The RG
logic was extended to include time-dependent forces on spatially fixed nuclei
directly \cite{KD86,Raj94A,RajBuot94B}. However, coupled electronic and
nuclear motion usually is treated via Li and Tong's extension \cite{LiTong86}
to multi-species systems. Some relevant issues regarding the Li-Tong
formulation are summarized in Appendix A. There is a considerable variety of
other work on the RG logic that, while of general relevance, is not specific
to the issues treated here \cite{Qian98,Chermyak00}.

For explicitly wave-function-based theory, a general approach to
non-Born-Oppenheimer dynamics has been given by Deumens, \"Ohrn, et al.
\cite{DeumensA,DeumensB}. They solve the time-dependent Schr\"odinger equation
(TDSE) via parameterization in group coherent states and the time-dependent
variational principle (TDVP). Here we take up the non-BO problem via extension
of our DFT analogue \cite{BLWSBT99a,BLWSBT00a} of the Deumens-\"Ohrn scheme.
The procedure includes a generalization to parameterization of the TDSE by
spatially varying functions (as distinct from the time-dependent complex
scalars that they use). The functions are the physical space-spin densities
that label the density operators via equivalence classes.

This formulation of the constrained search procedure
\cite{Levy79,LevyPerdew85,Lieb} via density operators has at least two
advantages. A general constrained search must include both pure states and
ensembles corresponding to a specified physical density. For density operators
the distinction between states and ensembles is not a barrier. Second, the
search is linear in the density operators but nonlinear in the states. The
density operator formulation therefore provides a unified way to expose the
internal structure of constrained search in forms exploitable for both
analysis and approximation. See anticipatory remarks along these lines in
Refs. \cite{Lieb} (see p.~263) and \cite{KD86} (just after eq. (5)).

Ref. \cite{BLWSBT99a} noted that our formulation implicitly applied to the
non-BO case. We develop that possibility in this and the following two papers.
This one gives the formal extension as a time-dependent constrained search.
Refs. \cite{BLWSBT99a,BLWSBT00a} are generalized from $N$-fermion operators
(and corresponding Hilbert spaces) alone to explicit inclusion of the nuclear
dynamics. Then we develop transformations that are constructive realizations
of constrained search. We give two realizations. Both use properties of the
contraction map from the $N$ -body density operator to the density. That map
confers a density-labeled fiber bundle structure on the density operators. The
kernel of the map plays a central role in parameterizing density evolution
paths. The more straightforward parameterization is by reference states in
each fiber plus linear combinations of elements of the kernel. It, however,
depends on a constraint which is difficult to implement. The more workable
parameterization of the dynamics comes from constructing the Lie groups that
map intra-fiber (i.e. at the same density) and those that map inter-fiber (one
density to another). Application of the TDVP provides equations of motion for
the density in terms of those Lie group parameters that, in the exact case, is
equivalent to the TDSE.

The second paper develops practical single-particle equations by restricting
the invariance transformations to one- and two-body contributions in order to
produce approximate time-dependent functionals by systematic refinement. The
third paper uses the invariance transformations to yield a path integral
version of TDDFT.

Because the development depends on considerable mathematical detail, this
paper is organized to present the principle results in the main text, with the
required supporting details in a set of Appendices.

\section{Density Mappings and Fiber Bundles}

\noindent This section shows how the densities at any time can be used to
classify density operators into a fiber bundle structure, then develops
exploitable properties of the bundle. This approach embodies constrained
search as a function of time in well-characterized, explicit transformations
along and among fibers. Systematic approximation of those transformations
provides a constructive route to approximate functionals (see paper II).

Without loss of generality we simplify to a single nuclear species (except if
otherwise noted). Let the generic electron and nuclear coordinates be
$\mathbf{r}$ and $\mathbf{R}$ respectively and similarly $s$, $S$ for the
spins. Composite coordinates are
\begin{equation}
\mathbf{x} \equiv( \mathbf{r} \, , \, s ) \qquad\mathbf{X} \equiv( \mathbf{R}
\, , \, S ) \qquad\mathbf{z} \equiv( \mathbf{x} \, , \, \mathbf{X} )
\end{equation}
and coordinate (space and spin) manifolds are
\begin{align}
\mathsf{x}  &  = \{ \mathbf{x}_{1},\mathbf{x}_{2}, \ldots\mathbf{x}_{N_{e}}
\}\nonumber\\
\mathsf{X}  &  = \{ \mathbf{X}_{1},\mathbf{X}_{2}, \ldots\mathbf{X}_{N_{n}} \}
\label{manifdef}%
\end{align}


To treat the electronic and nuclear degrees of freedom in balanced fashion and
thereby avoid (or at least minimize) problems inherent from decoupling
electron space and spin coordinates from the outset (e.g. symmetry breakings,
spin contamination \cite{PerdewSavin,BLWSBT98,SavinSem2}), requires the full
space-spin electron and nuclear density
\begin{equation}
\gamma(\mathbf{x},\mathbf{X},t)\equiv\gamma(\mathbf{z},t) \label{gammadef}%
\end{equation}
rather than the more common spatial density $n(\mathbf{r},\mathbf{\ R},t)$ or
spin-polarized spatial densities \cite{spinpol}. Note that $\gamma$ integrates
to the product $N_{e}N_{n}$, not their sum. $\gamma(\mathbf{x},\mathbf{X},t)$
can be interpreted as the joint probability for the simultaneous occurrence at
$t$ of an electron at position $\mathbf{x}$ and a nucleus at position
$\mathbf{X}$. This choice has the further benefit of avoiding possible
problems introduced by the translational invariance of isolated systems
\cite{KreibGross}. The single species space-spin densities can be recovered by
integration
\begin{align}
\gamma_{e}(\mathbf{x},t)  &  =\frac{1}{N_{n}}\int\gamma(\mathbf{x}%
,\mathbf{X},t)d\,\mathbf{X}\nonumber\\
\gamma_{n}(\mathbf{X},t)  &  =\frac{1}{N_{e}}\int\gamma(\mathbf{x}%
,\mathbf{X},t)d\,\mathbf{x}%
\end{align}


With two species, it is easiest to formulate the problem in second
quantization. The field operators are
\begin{align}
\Phi^{\dagger}_{e}(\mathbf{x})  &  = \sum_{i} \phi^{*}_{i}(\mathbf{x}%
)a_{i}^{\dagger}\nonumber\\
\Phi^{\dagger}_{n}(\mathbf{X})  &  = \sum_{j} \phi_{n,j}^{*}(\mathbf{X}%
)b_{j}^{\dagger} \label{fieldop}%
\end{align}
where the $\phi_{i}(\mathbf{x})$, $\phi_{n,j}(\mathbf{X})$ are elements of
complete one-particle basis sets for electrons and nuclei respectively and
$a_{i}^{\dagger}$, $b_{j}^{\dagger}$ are electron and nuclear creation
operators respectively, with commutation relations
\begin{align}
\left[  a_{i}^{\dagger}, a_{j} \right]  _{+}  &  = \delta_{ij} \qquad\left[
a_{i}^{\dagger}, a^{\dagger}_{j} \right]  _{+} = \left[  a_{i}, a_{j} \right]
_{+} = 0\nonumber\\
\left[  b_{i}^{\dagger}, b_{j} \right]  _{\pm}  &  = \delta_{ij} \qquad\left[
b_{i}^{\dagger}, b^{\dagger}_{j} \right]  _{\pm} = \left[  b_{i},b_{j}
\right]  _{\pm} = 0
\end{align}
Commutation or anti-commutation depends on the nuclear species; all
inter-species commutators are zero. Products of these operators for both an
electron and a nucleus will be written in compressed form for such products
\begin{equation}
\Phi^{\dagger}(\mathbf{z})\Phi(\mathbf{z}) \equiv\Phi^{\dagger}_{e}%
(\mathbf{x}) \Phi_{e}(\mathbf{x})\Phi^{\dagger}_{n}(\mathbf{X}) \Phi
_{n}(\mathbf{X}) \label{fieldopprod}%
\end{equation}


Generalizing Refs.~\cite{BLWSBT99a,BLWSBT00a}, the total system density
operator $D^{N}(t)$ is defined in the tensor product Hilbert space of the
$N_{e}$ electrons and $N_{n}$ nuclei, $N=N_{e}+N_{n}$. The time-dependent
density then is
\begin{equation}
\gamma(\mathbf{z},t)=Tr\left\{  \Phi^{\dagger}(\mathbf{z})\Phi(\mathbf{z}%
)D^{N}(t)\right\}  =\gamma(\mathbf{x},\mathbf{X},t)
\end{equation}
In the coordinate representation, the total system wave function
$\Psi(\mathsf{x},\mathsf{X},t)$ expressed in terms of the vacuum state
$|0\rangle$ is
\begin{align}
\Psi(\mathsf{x},\mathsf{X},t)  &  =\langle\mathbf{x}_{1},\ldots,\mathbf{x}%
_{N_{e}},\mathbf{X}_{1},\ldots,\mathbf{X}_{N_{n}}|\Psi(t)\rangle\nonumber\\
&  =\langle0|\Phi_{e}(\mathbf{x}_{1})\ldots\Phi_{e}(\mathbf{x}_{N_{e}}%
)\Phi_{n}(\mathbf{X}_{1})\ldots\Phi_{n}(\mathbf{X}_{N_{n}})\Psi(t)\rangle
\end{align}
where $|\Psi(t)\rangle\in\mathcal{H}^{N_{e}}\otimes\mathcal{H}^{N_{n}}$ and
the associated $N$-body density operator integral kernels are
\begin{align}
&  D^{N}(\mathbf{x}_{1}^{\prime},\ldots\mathbf{x}_{N_{e}}^{\prime}%
,\mathbf{X}_{1}^{\prime},\ldots\mathbf{X}_{N_{n}}^{\prime};\mathbf{x}%
_{1},\ldots\mathbf{x}_{N_{e}},\mathbf{X}_{1},\ldots\mathbf{X}_{N_{n}%
},t)\nonumber\\
&  =\Psi^{\ast}(\mathbf{x}_{1}^{\prime},\ldots\mathbf{x}_{N_{e}}^{\prime
},\mathbf{X}_{1}^{\prime},\ldots\mathbf{X}_{N_{n}}^{\prime},t)\Psi
(\mathbf{x}_{1},\ldots\mathbf{x}_{N_{e}},\mathbf{X}_{1},\ldots\mathbf{X}%
_{N_{n}},t)
\end{align}
(Remark that $D^{N}$ are called density matrices in the theoretical chemistry
literature, which is inconvenient terminology here.) The one-electron density
operator kernel $D_{e}^{1}(\mathbf{x}^{\prime},\mathbf{x},t)$ arises from
integrating out all the other electron coordinates \emph{and all} the nuclear
coordinates. The converse holds for the nuclear density operator kernel. In
the BO approximation, the electron density operator kernel is a conditional
probability relative to the nuclear positions. A relevant discussion is in
Ref. \cite{Sjoqvist}.

With one substantive alteration, the formulation in
Refs.~\cite{BLWSBT99a,BLWSBT00a}, involving trace class operators,
Hilbert-Schmidt operators, and contraction maps to one-particle function
spaces goes through essentially unchanged except for obvious nuclear spin
symmetry details. The alteration is a switch to unnormalized states. These are
convenient in the TDVP, where explicit normalization is used.

The central quantity for constrained search is the contraction map that yields
the combined density from the density operator
\begin{align}
\Xi &  :D^{N}(t)\rightarrow\gamma(\mathbf{z},t)\nonumber\\
\gamma(\mathbf{z},t)  &  =Tr\{\Phi^{\dagger}(\mathbf{z})\Phi(\mathbf{z}%
)D^{N}(t)\} \label{Xidefn}%
\end{align}
Here
\begin{align}
Tr\{D^{N}(t)\}  &  =1\;\;\forall\;t\nonumber\\
\int\gamma(\mathbf{z},t)d\mathbf{z}  &  =N_{e}N_{n}\;\forall\;t
\end{align}
$\Xi$ itself is a linear map from the space $\mathcal{B}_{1}(\mathcal{H}^{N})$
of trace class operators on $N$-particle Hilbert space to absolute integrable
functions on $\mathbb{R}^{6}\times\mathbb{C}^{4}$.

The density labels equivalence classes $[\gamma]_{N}$ of positive trace-class
operators $Y$ and $W$ by
\begin{equation}
Y\,,W\in\lbrack\gamma]_{N}\Leftrightarrow\Xi(Y)=\Xi(W)=\gamma
\end{equation}
These equivalence classes are elements of a bundle \cite{Bruhat}. By
inspection it is a fiber bundle (specifically an affine bundle)
\cite{affinebundleref}; see below. Denote it as $\mathfrak{B}_{aff}$. The base
of the bundle is the set of densities, which belong to $L_{1}(\mathbb{R}%
^{6}\times\mathbb{C}^{4})$. The elements of one fiber are the members of an
equivalence class which form an affine space; see below. The identification follows.

If restricted to normalized, positive, trace-class operators, this
classification is physically equivalent to the original constrained search
scheme, where states are grouped by the density they yield. The analytical
advantage is that the problem is linear in the density operators but nonlinear
in the states. To realize that advantage we extend consideration, temporarily,
to all trace class operators. Two consequences must be addressed eventually
(i) not all the quantities $f$ generated by the contraction map $\Xi$ are
physical densities $\gamma$ (everywhere positive), (ii) for a genuine physical
density $\gamma$ not every trace-class operator that contracts to it is
positive. With these caveats, the extended equivalence class is
\begin{equation}
\lbrack f]_{N}=\left\{  Y\in\mathcal{B}_{1}(\mathcal{H}^{N});\Xi(Y)=f\right\}
\label{genequiv}%
\end{equation}
(Aside different $f$'s do not necessarily have the same norm.) The gain is
that the kernel of $\Xi$
\begin{equation}
\ker\{\Xi\}=\{Y;Tr\,[\Phi^{\dagger}\Phi Y]=0\}\equiv\left\{  U-W;\Xi\left(
U\right)  =\Xi\left(  W\right)  \right\}  \label{kerdef}%
\end{equation}
is a linear subspace of $\mathcal{B}_{1}(\mathcal{H}^{N})$. Then every element
in a given equivalence class is expressible as
\begin{equation}
Y=\sum_{i=1}^{d}\alpha_{i}W_{i};\quad\sum_{i=1}^{d}\alpha_{i}=1
\label{YexpandW}%
\end{equation}
(which shows that $[f]_{N}$ is an affine space), and
\begin{equation}
W_{i}=Y_{ref}+\xi_{i};\;\;i>0 \label{WifromYref}%
\end{equation}
Here $Y_{ref}$ is a fixed, arbitrary element of the particular $[f]_{N}$ (i.e.
is not in $\ker\{\Xi\}$) and the $\{\xi_{i}\}$ constitute a basis for $\ker
\Xi$ \cite{Yref}. Each fiber thus is identified by a reference element in
$[f]_{N}$, with its internal structure characterized by the kernel via eqs.
(\ref{YexpandW}, \ref{WifromYref}). Eq. (\ref{YexpandW}) can be recast as
\begin{align}
Y  &  =Y_{ref}+U\nonumber\\
U  &  =\sum_{i=1}^{d}\alpha_{i}\xi_{i} \label{YYref}%
\end{align}
and the bundle is affine \cite{diffbundle}.

It follows that the orthogonal projector onto any pure $N$-particle state
$|\Psi_{N}\rangle$ with density $\gamma$ can be written as
\begin{equation}
\frac{\left|  \Psi_{N}(\mathbf{c})\right\rangle \left\langle \Psi
_{N}(\mathbf{c})\right|  }{\left\langle \Psi_{N}(\mathbf{c})|\Psi
_{N}(\mathbf{c})\right\rangle }=Y_{ref}(\gamma)+\sum_{i=1}^{d}c_{i}\xi_{i}
\label{Psicxi}%
\end{equation}
This expression is the essence of pure-state constrained search. Changes in
values of the $\{c_{i}\}$ move within a fiber, i.e. search over all system
density operators corresponding to a fixed density. Changing to another fiber
means another density has been selected. Both TDDFT and ground state DFT
result by combination with the appropriate variation principle (stationary
action for TDDFT, Rayleigh-Ritz for ground state).

\section{Explicit Implementation of Time-dependent Constrained Search}

\noindent The next task is to implement the search as the solution of the
time-dependent Schr\"odinger equation (TDSE). Refs.~\cite{BLWSBT99a,BLWSBT00a}
extended the work of Kramer and Sarceno \cite{KramSar} (``K\&S'' versus ``KS''
for Kohn and Sham) to include parameterization of the TDSE via the density as
a function of time and electronic coordinate $\mathbf{x}$. This section gives
the extension for the combined density $\gamma(\mathbf{z, t)}$ and the
resulting dynamics in terms of the coefficients $\mathbf{c}$ in eq.
(\ref{Psicxi}). To exploit the fiber bundle properties the argument here
differs somewhat from Ref. \cite{BLWSBT99a}.

The original K\&S analysis of the TDSE can be extended straightforwardly to
the manifold of states $\left|  \Psi(\lambda(t))\right\rangle $ parameterized
by real, time-dependent functions $\lambda(t,\mathbf{z})$. For these states,
solutions of the TDSE (denoted $\left|  \Psi(\lambda_{\#}(t))\right\rangle $)
are characterized by parameterizing functions that satisfy
\begin{equation}
\int\,\eta(t,\,\mathbf{y},\,\mathbf{y}^{\prime})\frac{d\lambda(\mathbf{y}%
^{\prime})}{dt}\,d\mathbf{y}^{\prime}=\frac{\delta\mathcal{E}(t)}%
{\delta\lambda(\mathbf{y})}\qquad\qquad\forall\;\mathbf{y} \label{dEdlambda}%
\end{equation}
where
\begin{equation}
\eta(t,\,\mathbf{y},\,\mathbf{y}^{\prime})=i\left\{  \frac{\delta}%
{\delta\,\tilde{\lambda}(\mathbf{y})}\frac{\delta}{\delta\,\lambda
(\mathbf{y}^{\prime})}\,-\,\frac{\delta}{\delta\,\lambda(\mathbf{y})}%
\frac{\delta}{\delta\,\tilde{\lambda}(\mathbf{y}^{\prime})}\right\}
\ln\left\langle \Psi(\tilde{\lambda}(t))|\Psi(\lambda(t))\right\rangle
\,|_{\tilde{\lambda}=\lambda} \label{eta}%
\end{equation}
and
\begin{equation}
\mathcal{E}(t)=\frac{\left\langle \Psi(\lambda(t))|H|\Psi(\lambda
(t))\right\rangle }{\left\langle \Psi(\lambda(t))|\Psi(\lambda
(t))\right\rangle } \label{Eoft}%
\end{equation}
Note the explicit norm. The parameterizing function space metric satisfies
\begin{equation}
\eta(t,\,\mathbf{y},\,\mathbf{y}^{\prime})=-\eta(t,\,\mathbf{y}^{\prime
},\,\mathbf{y})
\end{equation}


Eq. (\ref{dEdlambda}) can be cast in a form analogous with classical
mechanics, to wit
\begin{align}
\frac{d\lambda(\mathbf{y}^{\prime})}{dt}  &  =\int\,\eta(t,\,\mathbf{y}%
^{\prime},\,\mathbf{y})^{-1}\frac{\delta\mathcal{E}(t)}{\delta\lambda
(\mathbf{y})}\,d\mathbf{y}\quad\forall\;\mathbf{y}^{\prime}\nonumber\\
\dot{\lambda}(\mathbf{y}^{\prime})  &  =\{\lambda(\mathbf{y}),\mathcal{E}\}
\label{genEOM}%
\end{align}
with the generalized Poisson bracket
\begin{equation}
\{f(t),g(t)\}\equiv\int\frac{\delta f(t)}{\delta\,\lambda(\mathbf{y})}%
\zeta(t,\,\mathbf{y},\,\mathbf{y}^{\prime})\frac{\delta g(t)}{\delta
\,\lambda(\mathbf{y}^{\prime})}d\mathbf{y}\,d\mathbf{y}^{\prime}
\label{genPoisson}%
\end{equation}
and the definition
\begin{equation}
\zeta(t,\,\mathbf{y},\,\mathbf{y}^{\prime})\equiv\eta^{-1}(t,\,\mathbf{y}%
,\,\mathbf{y}^{\prime}) \label{etainverse}%
\end{equation}
assuming that $\eta^{-1}$ exists. This is not a strongly restrictive
assumption \cite{etainv}.

\emph{If} the manifold of states $\left|  \Psi(\lambda(t))\right\rangle $ is
all of $\mathcal{H}^{N}$, \emph{then} solutions of Eqs. (\ref{genEOM}),
(\ref{genPoisson}) are equivalent to solutions of the full TDSE except for an
overall phase that also can be determined within the parameterization
\cite{KramSar,BLW94}. It follows from eq. (\ref{Psicxi}) that, indeed
\begin{equation}
\mathcal{H}^{N}=\left\{  \left|  \Psi_{N}(\mathbf{c},\gamma)\right\rangle
\right\}
\end{equation}
with an obvious change in notation to account for the $\gamma$ dependence.
Note that the set of coefficients $\mathbf{c}$ must be constrained such that
the result in eq. (\ref{Psicxi}) is a pure-state projector, in symbols
\begin{equation}
\{c_{i}\}\subset\Lambda_{Y_{ref}(\gamma)}\ni Y=\mathrm{pure \, state\, density
\,operator} \label{posconstraint}%
\end{equation}
Characterization of the set $\Lambda$ is deferred. The equation of motion
(\ref{genEOM}) that is fully equivalent to the TDSE becomes
\begin{equation}
\{(\mathbf{c,}\gamma),\mathcal{E}\}=(\dot{\mathbf{c,}}\dot{\gamma})
\label{cgammaeom}%
\end{equation}
Since all the quantities involved are functionals of $\gamma$, this equation
demonstrates the existence of a TDDFT independent of the RG logic. This claim
becomes evident upon partitioning eq. (\ref{cgammaeom}) between the
$\mathbf{c}(t)$ and $\gamma(t)$; see Appendix B. The outcome is that the
density dynamics can be displayed explicitly, at least formally, with the
$c_{i}$ as functionals of $\gamma$. In terms of them, the exact temporal path
for the density is the solution of
\begin{equation}
\dot{\gamma}(t)=\tilde{\mathcal{N}_{1}}(\gamma(t))\left[  \nabla_{\gamma
}\,\mathcal{E}(\gamma(t))-\tilde{\mathcal{N}_{2}}(\gamma(t))\right]
\label{gammadot1}%
\end{equation}
where $\tilde{\mathcal{N}_{1}}(\gamma(t))$, $\tilde{\mathcal{N}_{2}}%
(\gamma(t))$ are density functionals determined as self-consistent solutions
of the equation of motion for the coefficients $\mathbf{c}$. See Appendix B.

\section{Dynamics Parameterized by Group Properties of the Fiber Bundle}

Parameterization via the $\mathbf{c}$ coefficients has a clear linkage to the
structure of constrained search. There is no straightforward way, however, to
implement the constraint eq. (\ref{posconstraint}) requiring the $c_{i}$ to
preserve positivity of the operator $Y$ in eq. (\ref{YYref}). Observe that
this constraint is in addition to having to select an approximation to model
the states $\left|  \Psi_{N}(\mathbf{c},\gamma)\right\rangle $. A suggestive
aspect of the $\mathbf{c}$ parameterization is that in general it is not
complex analytic. A complex analytic parameterization is technically useful
for dealing with the positivity constraint, as shown next. Such a
parameterization with independent real and imaginary parts always is
equivalent to a real parameterization, hence does not violate the assumptions
that give eqs. (\ref{dEdlambda}) - (\ref{Eoft}).

Thus we develop an alternative parameterization based on group properties of
transformations of a vector bundle, $\mathfrak{B}_{vect}$, that is isomorphic
to the affine bundle, $\mathfrak{B}_{aff}$, constructed in the previous
Section. Use of this vector bundle does lead to a complex analytic description
and also makes the group theoretic analysis less demanding (we only need to
consider transformations of vector spaces rather than affine spaces).
$\mathfrak{B}_{vect}$ is based on the direct sum decomposition of
$\mathcal{B}_{1}\left(  \mathcal{H}^{N}\right)  $ given by
\begin{equation}
\mathcal{B}_{1}\left(  \mathcal{H}^{N}\right)  =\ker\Xi\oplus\ker\Xi^{c}
\label{directsum}%
\end{equation}
It has a base space $\ker\Xi^{c}$, the vector space complement of $\ker\Xi$ in
$\mathcal{B}_{1}\left(  \mathcal{H}^{N}\right)  $, which is isomorphic to
$L_{1}(\mathbb{R}^{6}\times\mathbb{C}^{4})$, and fibers $\ker\Xi$ that are all
identical. Note therefore that the base of $\mathfrak{B}_{vect}$ is isomorphic
with the set of absolutely integrable functions that contain the set of
densities and that each equivalence class $[f]_{N}$ is comprised of the fiber
plus a particular base element. Both the fibers and the base are vector spaces.

Transformations of $\mathcal{B}_{1}\left(  \mathcal{H}^{N}\right)  $ that are
block-triangular with respect to bases adapted to this decomposition, that is,
transformations of the form
\begin{equation}
\left(
\begin{array}
[c]{cc}%
\left|  \ker\Xi\right)  \left(  \ker\Xi\right|  & \left|  \ker\Xi\right)
\left(  \ker\Xi^{c}\right| \\
0 & \left|  \ker\Xi^{c}\right)  \left(  \ker\Xi^{c}\right|
\end{array}
\right)  \label{blockxform}%
\end{equation}
induce vector bundle maps $\mathfrak{B}_{vect}\rightarrow\mathfrak{B}_{vect}$
and form the group, $\mathfrak{G}_{\Xi}$, of automorphisms of $\mathfrak{B}%
_{vect}$ \cite{Bruhat}. Here $\left|  A\right)  \left(  B\right|  $ denotes
mappings from the subspace $B$ to the subspace $A$ and the lower block of eq.
(\ref{blockxform}) is zero because mappings from the fiber to the base would
destroy the vector bundle properties.

Usually invariance groups characterizing a property of a class of states are
constructed from unitary transformations. In the present case there is no
unitary group of transformations that operates homogeneously, (i.e. produces
all pure states from any given reference pure state), both in $\mathcal{H}%
^{N}$ and the set of pure states contained in each equivalence class $\left[
\gamma\right]  _{N}$, and also preserves the vector bundle structure of
$\mathfrak{B}_{vect}$. However one can construct a commutative, but
non-unitary, Lie group, $\mathfrak{G}_{\ker}$, of transformations contained in
the general linear group, $GL\left(  \mathcal{H}^{N}\right)  $, of
$\mathcal{H}^{N}$ that does have these properties. This group induces positive
congruence maps of $\mathcal{B}_{1}\left(  \mathcal{H}^{N}\right)
\mapsto\mathcal{B}_{1}\left(  \mathcal{H}^{N}\right)  $, hence maps states to
states (details in Appendix C). The group $\mathfrak{G}_{\ker}$ contains
subgroups $\mathfrak{G}_{den}$ and $\mathfrak{G}_{fib}$ that can be used to
give parameterizations of $\ker\Xi^{c}$ and $\ker\Xi$ respectively.

The parameterization of the base, $\ker\Xi^{c}$, that results is non-
redundant (i.e. one set of parameters corresponds to one and only one element)
and is given in terms of the group $\mathfrak{G}_{den}\subset\mathfrak{G}%
_{\ker}$ which is represented by transformations of the form
\begin{equation}
\left(
\begin{array}
[c]{cc}%
t\left(  T_{B}\right)  & 0\\
0 & T_{B}%
\end{array}
\right)
\end{equation}
acting on a reference element in $\ker\Xi^{c}$, where the top block is
determined by the bottom one. In the case of the fibers, $\ker\Xi$, we produce
a non-redundant parameterization of \textit{pure states }given in terms of the
group $\mathfrak{G}_{fib}$ of transformations
\begin{equation}
\left(
\begin{array}
[c]{cc}%
T_{F} & 0\\
0 & I_{B}%
\end{array}
\right)
\end{equation}
acting on a pure state reference, where the bottom block is the identity
operation on the base. SBTSBT next sentence not specific; which Appendix or
Appendices? SBTSBT A more detailed discussion of the properties of these
transformations, which are determined by integral kernels of operators acting
in $\mathcal{H}^{N}$, that are diagonal in the coordinate representation, are
given later in this section and in the Appendices.

States (pure or ensemble) of a physical system correspond to positive trace
class operators. For pure states these operators correspond to elements of
$\ \mathcal{H}^{N}$. As we are interested in mappings between states, positive
maps of such operators are critical. Since all physical densities can be
associated with pure states, we focus on them. Relevant details of positive
operators and maps are summarized in Appendix D. The outcome is that the only
two possibilities for mappings between pure states are unitary transformations
and the more general congruence transformations\cite{congruence} of
$\mathcal{B}_{1}\left(  \mathcal{H}^{N}\right)  $. To incorporate both types,
we consider the most general possible invertible transformations on
$\mathcal{H}^{N}$, namely those of the General Linear group $GL(\mathcal{H}%
^{N})$
\begin{equation}
T\,\in\,GL(\mathcal{H}^{N})\ni T|\Psi\rangle=|\chi\rangle\text{ and }%
T^{-1}|\chi\rangle=|\Psi\rangle\quad\forall|\Psi\rangle,|\chi\rangle
\in\mathcal{H}^{N}%
\end{equation}
Because norms appear explicitly in the TDVP, the non-unitarity of these maps
is not a difficulty. It follows that
\begin{equation}
T=e^{A}\,e^{iB} \label{eAeiB}%
\end{equation}
with $A^{\dagger}=A$ and $B^{\dagger}=B$ (see Appendix D). All pure state
density operators are mapped to one another by such transformations; see
Appendix E on this point along with the basis of a systematic extension of the
present analysis to ensemble density operators.

The most general group, $\mathfrak{G}_{\ker}\subset\mathfrak{G}_{\Xi}\subset
GL\left(  \mathcal{B}_{1}\left(  \mathcal{H}^{N}\right)  \right)  $, of
positive transformations of $\mathcal{B}_{1}\left(  \mathcal{H}^{N}\right)  $
that both preserve the fiber bundle structure of $\mathfrak{B}_{vect}$ and
also are congruence transformations then has the following property
\begin{equation}
\mathfrak{G}_{\ker}=\left\{  T\in GL(\mathcal{H}^{N});Tr\left\{  \Phi
^{\dagger}\Phi TYT^{\dagger}\right\}  =0\;\forall Y\,\in\ker\Xi\text{
}\right\}  \label{GLkerXi}%
\end{equation}
By resort to the coordinate eigenfunction representation, we show in
Appendices F and G that only certain subgroups of the commutative subgroup
$\mathcal{C}(\mathcal{H}^{N})$ of diagonal coordinate multi-particle
transformations $T(\Upsilon)$ have this property. In terms of the field
operators $\left\{  \Phi^{\dagger}\left(  \kappa\right)  ,\Phi\left(
\kappa\right)  ;\kappa=\left(  \mathbf{x},\mathbf{X}\right)  \right\}  $,
these diagonal transformations are
\begin{align}
\mathcal{C}(\mathcal{H}^{N})  &  =\left\{  T(\Upsilon);T(\Upsilon)YT^{\dagger
}(\Upsilon)\equiv\widehat{T}(\Upsilon)(Y)\right.  \,;\nonumber\\
T(\Upsilon)  &  =\left.  e^{\int\Upsilon(\kappa_{1},\ldots,\kappa_{N}%
)\Phi^{\dagger}(\kappa_{1})\Phi(\kappa_{1})\ldots\Phi^{\dagger}(\kappa
_{N})\Phi(\kappa_{N})}d\kappa_{1}\ldots d\kappa_{N}\right\} \nonumber\\
&  \Upsilon(\kappa_{1},\ldots,\kappa_{N})\in\mathbb{C}%
\end{align}
Here and in the Appendices the coordinate strings are in the compressed
notation
\begin{equation}
\kappa_{1}\ldots\kappa_{N_{e}},\kappa_{N_{e}+1}\ldots\kappa_{N}\equiv
\mathbf{x}_{1}\ldots\mathbf{x}_{N_{e}},\mathbf{X}_{1}\ldots\mathbf{X}_{N_{n}}
\label{compressed}%
\end{equation}
By construction $\mathcal{C}(\mathcal{H}^{N})$ has the property
\begin{equation}
\left[  T(\Upsilon),\Phi^{\dagger}\Phi\right]  =0
\end{equation}
Note that \emph{all} states (unnormalized in general) belonging to
$\mathcal{H}^{N}$ can be expressed as
\begin{equation}
|\Psi(\Upsilon)\rangle=T(\Upsilon)|\Psi\rangle
\end{equation}
for any fixed reference state $|\Psi\rangle\in\mathcal{H}^{N}$; see Appendix
H. Separation of $\Upsilon$ into real and imaginary parts as in eq.
(\ref{eAeiB}) gives
\begin{align}
T(\Upsilon)  &  =e^{\int\mathcal{V}(\kappa_{1},\ldots,\kappa_{N})\Phi
^{\dagger}(\kappa_{1})\Phi(\kappa_{1})\ldots\Phi^{\dagger}(\kappa_{N}%
)\Phi(\kappa_{N})d\kappa_{1}\ldots d\kappa_{N}}\nonumber\\
&  \times e^{i\int\omega(\kappa_{1},\ldots,\kappa_{N})\Phi^{\dagger}%
(\kappa_{1})\Phi(\kappa_{1})\ldots\Phi^{\dagger}(\kappa_{N})\Phi(\kappa
_{N})d\kappa_{1}\ldots d\kappa_{N}}\nonumber\\
&  \equiv e^{\breve{\mathcal{V}}}e^{i\breve{\omega}} \label{TVomega}%
\end{align}


The group of transformations $\mathcal{G}_{\Xi}$ generated by the imaginary
part of $T$ is a unitary subgroup of $\mathcal{C}(\mathcal{H}^{N})$ that
leaves $\ker\Xi$ invariant
\begin{align}
\!\mathcal{G}_{\Xi}  &  =\left\{  U(\omega);\hat{U}(\omega)(Y)\equiv
U(\omega)YU^{\dagger}(\omega);\right.  \omega(\kappa_{1},\ldots,\kappa_{N}%
)\in\mathbb{R}\nonumber\\
U(\omega)  &  =\left.  e^{i\,\int\omega(\kappa_{1},\ldots,\kappa_{N}%
)\Phi^{\dagger}(\kappa_{1})\Phi(\kappa_{1})\ldots\Phi^{\dagger}(\kappa
_{N})\Phi(\kappa_{N})d\kappa_{1}\ldots d\kappa_{N}}\right\}  \label{GXidefn1}%
\end{align}
Elements of this subgroup map each fiber to itself i.e. are intra-fiber maps
\begin{equation}
Tr\{\Phi^{\dagger}\Phi U(\omega)YU^{\dagger}(\omega)\}=Tr\{\Phi^{\dagger}\Phi
Y\}\quad\forall\,Y\in\mathcal{B}_{1}(\mathcal{H}^{N})
\end{equation}
In fact $\mathcal{G}_{\Xi}$ is a subgroup of \emph{the} invariance group of
the contraction map $\Xi$, since
\begin{equation}
\Xi\left(  U(\omega)YU^{\dagger}(\omega)\right)  =\Xi\left(  Y\right)  \text{
}\forall Y\in\mathcal{B}_{1}\left(  \mathcal{H}^{N}\right)
\end{equation}
The group $\mathfrak{G}_{\ker}$ can be expressed as the direct product
$\mathcal{G}_{\ker}\times\mathcal{G}_{\Xi}$, where the subgroup,
$\mathcal{G}_{\ker}$, is generated by $\mathcal{V}$ (the real part of
$\Upsilon$).

To construct a representation of this subgroup, which contains both inter- and
intra-fiber maps, restrictions must be imposed upon $\mathcal{V}(\kappa
_{1},\ldots,,\kappa_{N})$ so that eq. $\left(  \ref{GLkerXi}\right)  $ is
satisfied. In Appendix G we show that the diagonal integral kernel $T_{d}$ of
this transformation may be written as
\begin{equation}
|T_{d}(\kappa_{1},\ldots\kappa_{N})|^{2}=1+\phi(\kappa_{1},\ldots\kappa
_{N})=e^{2\mathcal{V}(\kappa_{1},\ldots,,\kappa_{N})}%
\end{equation}
(note Appendix G, eq. (\ref{phidef})), i.e., $T_{d}=e^{\breve{\mathcal{V}}%
}=\left(  1+\breve{\phi}\right)  ^{\frac{1}{2}}$. This decomposition is
significant because of the requirement (Appendix G) that $\breve{\phi}$ needs
to be strongly orthogonal to $K_{D}$, the diagonal part of $\ker\Xi$ in the
coordinate representation. That is, the integral kernel $\phi(\kappa
_{1},\ldots\kappa_{N})$ of $\breve{\phi}$ must be strongly orthogonal to the
integral kernels of all diagonal operators that lie in $\ker\Xi$. This
condition is obtained, (Appendix G), by considering the decomposition of
$Y\in\mathcal{B}_{1}(\mathcal{H}^{N})$ as
\begin{equation}
Y=Y_{K_{OD}}+Y_{K_{D}}+\Gamma\left(  \gamma\right)
\end{equation}
where we decompose $\ker\Xi$ into off diagonal and diagonal parts, i.e.
$\ker\Xi$ $=K_{OD}\oplus K_{D}$, and $Y_{K_{OD}}\in K_{OD}$, $Y_{K_{D}}\in
K_{D}$, and $\Gamma\left(  \gamma\right)  \in\ker\Xi^{c}$. Observe that the
transformed diagonal part $\left(  1+\breve{\phi}\right)  ^{\frac{1}{2}%
}Y_{K_{D}}\left(  1+\breve{\phi}\right)  ^{\frac{1}{2}}=\left(  1+\breve{\phi
}\right)  Y_{K_{D}}$ only belongs to $Y_{K_{D}}$ if $\phi\left(  \kappa
_{1},\ldots\kappa_{N}\right)  $ is strongly orthogonal to $Y_{K_{D}}\left(
\kappa_{1},\ldots\kappa_{N}\right)  $, while $\left(  1+\breve{\phi}\right)
^{\frac{1}{2}}Y_{K_{OD}}\left(  1+\breve{\phi}\right)  ^{\frac{1}{2}\dagger}$
always belongs to $K_{OD}$ (note that $\left(  1+\breve{\phi}\right)
^{\frac{1}{2}}$ does not commute with $Y_{K_{OD}}$) and $\left(  1+\breve
{\phi}\right)  \Gamma\left(  \gamma\right)  $ belongs to $\ker\Xi^{c\text{ }}$
both in the case when $\breve{\phi}$ is strongly orthogonal to $\Gamma\left(
\gamma\right)  $ and when it is not. Clearly this map maintains the vector
bundle structure of $\mathfrak{B}_{vect}$ as it does not send elements of the
fibers to the base $\ker\Xi^{c}$.

Not all elements of $\mathcal{G}_{ker}$ are intra-fiber maps. In general
$\left(  1+\breve{\phi}\right)  \Gamma\left(  \gamma\right)  =\Gamma\left(
\gamma^{\prime}\right)  +Y_{K_{D}}^{\prime}$, and thus $\left(  1+\breve{\phi
}\right)  $ is an intra-fiber map only if $\gamma^{\prime}=\gamma$. In terms
of eq. (\ref{YYref}), the issue is that elements of $\mathcal{G}_{ker}$ may
change $Y_{ref}\rightarrow Y_{ref}^{\prime}$, where $Y_{ref}^{\prime}$ is in a
different fiber, to $Y_{ref}$. However, $\mathcal{G}_{ker}$ does contain a
subgroup, $\mathcal{G}_{fib\text{,}}$ of intra fiber maps; see Appendix G
also. An essential result is that the $\mathcal{\breve{V}}_{fib}$ that
generate this group are expressed in terms of $\breve{\phi}$ that are strongly
orthogonal to $K_{D}$ and act as the zero operator on $S_{D}$. Together with
the group $\mathcal{G}_{\Xi}$ the group $\mathcal{G}_{fib}$ forms, via direct
product, the group $\mathfrak{G}_{fib}$ of intra fiber super multiplication
operators, acting in $\mathcal{B}_{1}\left(  \mathcal{H}^{N}\right)  $ mapping
$\ker\Xi\mapsto\ker\Xi,$ generated by multiplication operators (in the
coordinate representation) acting in $\mathcal{H}^{N}$ i.e. $\mathfrak{G}%
_{fib}=\mathcal{G}_{fib}\times$ $\mathcal{G}_{\Xi}$. This group can be used to
give a non-redundant parameterization of $\ker\Xi$, as shown in Appendix H.

A representation of the group, $\mathfrak{G}_{den}$, of maps between fibers,
that is contained in $\mathcal{G}_{ker}$, can be constructed (Appendix G)
using $\mathcal{\breve{V}}_{den}$ based on $\breve{\phi}$ that are strongly
orthogonal to all diagonal operators that lie in $\ker\Xi$ but are not
strongly orthogonal to $S_{D}$. A non-redundant parameterization of $S_{D}$
based on a specific (but arbitrary) reference state $D_{ref}^{N}\left(
\gamma_{0}\right)  $ can be formed using this representation (Appendix G). It
is possible to choose the reference state $D_{ref}^{N}\left(  \gamma
_{0}\right)  $ and all the states generated from it, (one for each fiber), to
be Independent Particle States or to be states of any prescribed form.

Thus a non-redundant parameterization of the set of pure states in
$\mathfrak{B}_{vec}$ can be given in terms of the action of the product
$\mathfrak{G}_{fib}\times\mathfrak{G}_{den}=\mathcal{G}_{fib}\times
\mathcal{G}_{\Xi}\times\mathfrak{G}_{den}$ $\equiv\mathfrak{G}_{pure}$ on a
pure state reference element. Elements of $\mathcal{G}_{fib}\times
\mathcal{G}_{\Xi}$ are intra-fiber maps, while elements of $\mathfrak{G}%
_{den}$ are inter-fiber maps. The state $D_{ref}^{N}\left(  \mathcal{V}%
_{den}\right)  \equiv D_{ref}^{N}\left(  \mathcal{\gamma}\right)  $ obtained
through the non-redundant action of $\mathfrak{G}_{den}$ on $D_{ref}%
^{N}\left(  \gamma_{0}\right)  $ can be used as the reference state for the
non-redundant parameterization of the pure states contained in $\left[
\gamma\right]  _{N}$ through the action of $\mathfrak{G}_{fib}$ on
$D_{ref}^{N}\left(  \mathcal{\gamma}\right)  $, that gives $D^{N}\left(
\mathcal{V}_{fib}\left(  \gamma\right)  ,\mathcal{\omega}\right)  $. Note that
these maps are rank-preserving. The physical significance is that while these
maps do \emph{not} explore all elements of every fiber, they \emph{do} explore
all pure-state elements of every fiber if the initial reference is a pure
state. Technical details are in Appendices G and I.

The following diagram collects the primary relationships described in this
section. The top line is the domain and the bottom line the range for the
mappings shown in the middle.

%\begin{gather}
%\underbrace{{\frak B}_{aff}\backsimeq {\frak B}_{vect}\backsimeq {\cal B}%
%_{1}\left( {\cal H}^{N}\right) \equiv \ker \Xi \oplus \ker \Xi ^{c}}%
%\mathstrut \medskip  \nonumber \\
%\left\downarrow
%\begin{array}{c}
%{\frak G}_{\Xi }\subset GL\left( {\frak B}_{1}\left( {\cal H}^{N}\right)
%\right) \\
%\cup \\
%{\frak G}_{\ker }={\cal G}_{\ker }\times {\cal G}_{\Xi } \\
%\cup \\
%{\frak G}_{pure}={\frak G}_{fib}\times {\frak G}_{den} \\
%{\LARGE \parallel } \\
%{\cal G}_{fib}\times {\cal G}_{\Xi }\times {\frak G}_{den}%
%\end{array}%
%\right\downarrow  \nonumber \\
%\overbrace{{\frak B}_{aff}\backsimeq {\frak B}_{vect}\backsimeq {\cal B}%
%N_A{1}\left( {\cal H}^{N}\right) \equiv \ker \Xi \oplus \ker \Xi ^{c}}
%\nonumber \\
%{\frak G}_{\ker }\subset {\cal C}({\cal H}^{N})\cap \text{Congruence
%Transforms of }{\cal B}_{1}\left( {\cal H}^{N}\right)
%\end{gather}


\section{Equations of Motion}

As the final step, we use the parameterization of pure states in
$\mathcal{H}^{N}$ provided by the factorization of $\mathfrak{G}_{pure}$ in
terms of $\mathcal{G}_{\Xi}$, $\mathcal{G}_{fib}$, and $\mathcal{G}_{den}$ to
express the TDVP in terms of functionals of the physical density. We focus on
a form amenable to sequential approximation. First a reference state
$|\Psi_{0}\rangle\in\mathcal{H}^{N}$ must be selected \cite{RefState}. Given
$|\Psi_{0}\rangle$ and its associated density $\gamma_{0}$, \emph{any} pure
state (in general, unnormalized) in $\mathcal{H}^{N}$ can be parameterized in
terms of the foregoing group factorization as
\begin{equation}
|\Psi(\mathcal{V}_{fib}(\gamma),\omega,\mathcal{V}_{den}\rangle
=e^{_{\mathcal{\breve{V}}_{fib}(\gamma)}}\,e^{i\breve{\omega}}e^{\breve
{\mathcal{V}}_{den}}|\Psi_{0}\rangle\label{vv1omega1}%
\end{equation}
where
\begin{align}
e^{_{\mathcal{\breve{V}}_{fib}(\gamma)}}  &  \in\mathcal{G}_{fib}\nonumber\\
e^{\breve{\mathcal{V}}_{den}}  &  \in\mathfrak{G}_{den}\nonumber\\
e^{i\breve{\omega}}  &  \in\mathcal{G}_{\Xi} \label{vv1omega2}%
\end{align}
and $e^{\breve{\mathcal{V}}_{den}}$ produces states of different density,
$e^{_{\mathcal{\breve{V}}_{fib}(\gamma)}}$ produces states with the same
density as $e^{\breve{\mathcal{V}}_{den}}|\Psi_{0}\rangle$ and the shorthand
notation for exponentiated integral operators was introduced at
eq.~(\ref{TVomega}).

Equations (\ref{vv1omega1}), (\ref{vv1omega2}) make clear the distinction
between the time propagation of the density $\gamma(t)$ and the states. In
particular, it is vital that different $\breve{\mathcal{V}}_{den}$'s produce
states with different densities and that all possible densities $\gamma$
associated with pure states can be generated by varying $\mathcal{V}%
_{d}(\gamma_{0})$ such that there is a 1-1 correspondence between
$\breve{\mathcal{V}}_{den}$ and $\gamma$. That is, for each $\gamma$ there is
a unique transformation $e^{\breve{\mathcal{V}}_{den}}\in\mathfrak{G}_{den}$
that maps $|\Psi_{0}\rangle\mapsto|\Psi_{\gamma}\rangle$, where $|\Psi
_{\gamma}\rangle$ is \emph{some} state that yields $\gamma$. Details are in
Appendices E and H. Parameterization by $\breve{\mathcal{V}}_{den}$ therefore
is equivalent to parameterization with respect to $\gamma$. Equivalently
\begin{equation}
|\Psi(\mathcal{V}_{fib}(\gamma),\omega,\gamma)\rangle\equiv|\Psi
(\mathcal{V}_{fib}(\gamma),\omega,\mathcal{V}_{den}\rangle
\end{equation}
Notice in this last expression that the roles of $\mathcal{V}_{fib}(\gamma)$
and $\omega$ are to determine the time propagation of the states (hence also
the density operators) associated with $\gamma$. The ``fiber factor''
\begin{equation}
e^{_{\mathcal{\breve{V}}_{fib}(\gamma)}}\,e^{i\breve{\omega}}\equiv
e^{\breve{\Upsilon}} \label{fiberfactor}%
\end{equation}
selects among elements of the fiber, while $e^{\breve{\mathcal{V}}_{den}}$ is
the ``base factor'' that determines which $\gamma$ will be reached from the
initial one. From this perspective, the task of constructing density
functionals is equivalent to the task of finding the optimal $\mathcal{V}%
_{fib}(\gamma)$ and $\omega$ for the particular $\gamma$. By comparison with
the development leading to eq. (\ref{genEOM}), the equation of motion in terms
of this parameterization is
\begin{equation}
\{(\mathcal{V}_{fib}(\gamma),\omega,\gamma),\mathcal{E}\}=(\mathcal{\dot{V}%
}_{fib}(\gamma),\dot{\omega},\dot{\gamma}) \label{Vomegagammaeom}%
\end{equation}


Appendix J shows how to utilize the $\left(  \mathcal{V}_{fib}(\gamma
),\omega\right)  $ parameterization in a form which emphasizes the complex
analytic character of the fiber parameterization via use of the complex fiber
parameter
\begin{equation}
\Upsilon=\mathcal{V}_{fib}(\gamma)+i\omega\label{Ffactor}%
\end{equation}
The result of this parameterization is a different EOM, again a pure density
functional dynamical expression, of the same form as eq. (\ref{gammadot1})
\begin{equation}
\dot{\gamma}(t)=\mathcal{N}_{1}(\gamma(t))\left[  \nabla_{\gamma}%
\,\mathcal{E}(\gamma(t))-\mathcal{N}_{2}(\gamma(t))\right]  \label{gammadot2}%
\end{equation}
but with different functionals and with the positivity constraint built in.
Here $\mathcal{N}_{i}(\gamma(t))$, $1\leq i\leq2$, and $\mathcal{E}%
(\gamma(t))$ are density functionals determined by the solutions of the
equations of motion for the parameterizing functions $\mathcal{V}_{fib}%
(\gamma,t)$ and $\omega(t)$. Appendix J gives their detailed structure.

The $\mathcal{N}_{i}(\gamma(t))$ in eq. (\ref{gammadot2}) are fiber quantities
that provide information on intra-fiber searches and their propagation in
time. Thus they act as time-dependent and density-dependent potentials.
$\mathcal{E}(\gamma(t))$ also is an implicit function of fiber factors but an
explicit function of $\gamma$, hence may be decomposed in conventional KS
fashion into kinetic, coulombic, and exchange-correlation parts if that proves useful.

\section{Summary and Conclusions}

We summarize. First, to understand the details of constrained search in the
most general way, we established that the space of trace class operators
$\mathcal{B}_{1}(\mathcal{H}^{N})$ can be sorted into equivalence classes
\begin{equation}
\left\{  \lbrack f]_{N};\quad f\in L_{1}(\mathtt{Z})\right\}
\label{summequiv}%
\end{equation}
This classification endows $\mathcal{B}_{1}(\mathcal{H}^{N})$ with two fiber
bundle structures, one affine $\mathfrak{B}_{aff}$, the other vector
$\mathfrak{B}_{vect}$. In $\mathfrak{B}_{aff}$, the fibers are the $[f]_{N}$
and the base of the bundle is $L_{1}(\mathtt{Z})$, while in $\mathfrak{B}%
_{vect}$ the fibers are $\ker\Xi$ and the base is $\ker\Xi^{c}$. The
projection map from the fiber to the base in the affine case is the
contraction map $\Xi$ defined in eq. (\ref{Xidefn})
\begin{equation}
\Xi(Y)=Tr\{\Phi^{\dagger}\Phi Y\}=f;\quad Y\in\mathcal{B}_{1}(\mathcal{H}%
^{N});\quad f\in L_{1}(\mathtt{Z}) \label{summXidefn}%
\end{equation}
and the product of field operators is as in eq. (\ref{fieldopprod}). In the
vector bundle the corresponding projection map is the projection onto the
unique element $\Gamma\left(  f\right)  $ of $\ker\Xi^{c}$ that corresponds to
$f$ \ under the mapping $\Xi$. The fibers of the affine bundle (the
equivalence classes $[f]_{N}$ ) correspond to $\left\{  \ker\Xi\right\}
\times\left\{  \Gamma\left(  f\right)  \right\}  $ i.e. the set $\left\{
\left(  k,\Gamma\left(  f\right)  \right)  ;k\in\ker\Xi,\text{ fixed }%
\Gamma\left(  f\right)  \in\ker\Xi^{c}\right\}  $ which is a fiber and base
point of the vector bundle $\mathfrak{B}_{vect}$.

Second, because $\Xi$ is a linear map, its kernel is a linear subspace of
$\mathcal{B}_{1}(\mathcal{H}^{N})$. Because of the fiber bundle structures,
the kernel plays a central role in executing the part of constrained search
that corresponds to searching at a specified density. Detailed knowledge of
the structure of the kernel, in particular its invariances, therefore provides
an entirely new route to construction of functionals.

Third, because physical systems are characterized by density operators, i.e.
positive trace class operators $Y\in\mathcal{B}_{1}(\mathcal{H}^{N})_{+}%
\equiv\left\{  Y\in\mathcal{B}_{1}(\mathcal{H}^{N}),Y\geq0\right\}  $, we
focused upon fibers corresponding to $\left\{  [\gamma]_{N},\;\gamma
\geq0\right\}  $. Then we demonstrated [Appendix F, eq. (\ref{decomp})] that,
for pure states (and all physical densities correspond to pure states, though
not necessarily ground states), any density operator can be decomposed as
\begin{equation}
D^{N}=D_{K_{OD}}^{N}+D_{K_{D}}^{N}+\Gamma\left(  \gamma\right)
\end{equation}


Fourth, each of these three contributions has specific invariance properties
related to whether the operators are diagonal or off-diagonal in the
coordinate eigenfunction representation and to whether they lie in $\ker\Xi$
or not; see Appendices F and G. For time-dependent DFT, these invariances
provide a Lie group structure that can be used to parameterize the dynamics
via the TDVP. For both time-dependent and time-independent DFT, the
invariances provide a new set of tools for constructing functionals. We
distinguish ``constructing'' from ``devising'' because Appendices H and J
comprise a \emph{systematic} way to execute intra- and inter-fiber searches
via Lie groups of transformations. In both time-dependent and time-independent
forms of the theory, explicit knowledge of the optimal result of the
intra-fiber search is identical with determination of the exact functional for
the density which labels that particular fiber. When restricted to a
particular form of pure state, such knowledge is equivalent to determination
of the proper functional to associate with that type of state.

The analytic form of those Lie group transformations is suggestive of certain
well-known cluster expansions.\cite{Blaizot86} This aspect is taken up in Part
II (following paper) as one of several means to obtain workable systematic
approximations. Also treated there are the related issues of the various
classical approximations for the nuclear motion (clamped nucleus, adiabatic
motion on electronic potential, narrow wave-packet, etc.). On that topic we
remark here only that an effort to extract one of those results simply as the
limiting case of the exact dynamical equations must be treated with care
because such a limit requires imposition of a constraint on the Hamiltonian
(or action), hence amounts to invoking a different Hamiltonian to derive the
dynamical equations.

Finally, the challenge of constructing physically realistic, yet tractable
time-dependent functionals is evident from the intricate structure of the
$\mathcal{N}_{i} (\gamma(t))$ in eq. (\ref{gammadot2}). Further, because that
structure is related directly to the metric functions $\eta$, it is clear that
approximate functionals constructed on physical grounds or by generalization
from functionals used successfully in the time-independent case actually must
involve some implicit, unanalyzed assumptions about that metric.

\section{Acknowledgements}

We thank Evaert Jan Baerends, Kieron Burke, Erik Deumens, Yngve \"{O}hrn,
A.K.~Rajagopal, and Robert van Leeuwen for informative and helpful
discussions. SBT was supported in part by the U.S. National Science Foundation
(DMR-KDI 9980015).

\appendix


\section{Runge-Gross TDDFT for Electrons and Nuclei - Limitations}

\noindent Both the need for the Li-Tong multi-component scheme and its
limitations can be seen by considering the original one-component RG TDDFT.
For simplicity take $N_{n}$ nuclei of one species and $N_{e}$ electrons. RG
TDDFT assumes that the external dynamic potential
\begin{equation}
\hat V = \sum_{i}^{N_{e}} v(\mathbf{r}_{i},t) \label{VextofT}%
\end{equation}
has a temporal Taylor series expansion about the initial time $t=t_{0}$. Let
the initial spatial density be $n(\mathbf{r}, t=t_{0}) \equiv n_{0}%
(\mathbf{r})$. Then invertibility of the potential-density map is demonstrated
under this assumption by considering the contrary, two external potentials
$v^{\prime}(\mathbf{r},t)$, $v(\mathbf{r},t)$ that yield the same $n$ and
differ by more than a function of time alone. Nonexistence of such duplicate
potentials requires that a certain surface integral vanish. RG argued that it
vanishes for any physically meaningful potential, though they noted that the
Taylor series assumption disqualifies adiabatically switched-on potentials.
van Leeuwen provided an extension of that original argument, which, however,
does not address the issues raised below \cite{Leeuwen99}.

Examples are known for which the surface integral is not zero \cite{Xu85} but
they have been charcterized as non-physical\cite{Dhara87}. The latter
reference noted that the RG proof is not valid for both the homogeneous
electron gas \cite{HEGexample} and systems with moving bare nuclei. For the
latter, the temporal Taylor series has singularities of arbitrarily high order
\cite{bareTaylor}. The coefficients of the highly singular terms are time
derivatives of $X(t)$ evaluated at $t_{0}$. Those all vanish for stationary
nuclei and the most singular terms remaining in the Taylor series go as the
ordinary nuclear-electron potential. This problem is not discussed explicitly
in Ref. \cite{Saalmann96} [at their equation (12)] but see below.

To handle the surface integral, Gross and Kohn introduced two requirements,
one physical, the other formal \cite{GrossKohn90}. The formal constraint (on
the asymptotic behavior, in $r$, of the Taylor series coefficients of
$v(\mathbf{r}_{i},t)$) is not relevant here. The physical requirement is that
``experimentally realizable potentials are always produced by \emph{real}
(this means, in particular, normalizable) external charge densities
$\rho_{ext}(\mathbf{r}t)$ \ldots.'' Clearly this requirement does not handle
moving bare nuclei as the external charge densities, even though such
delta-function charges are normalizable in the usual sense. It might be argued
that non-singular distributions were intended. For nuclear motion this means
screening, in practice, a pseudopotential. This is what Refs.
\cite{Teil92,Fois88} do in the BO approximation. Refs.
\cite{Saalmann96,Knospe99} do the same in practice. Non-singular distributions
are fine practically, but do not solve the basic problem. Restriction to such
distributions would force TDDFT with nuclear motion never to be all-electron.
In turn that would lead to a basic inconsistency the time-independent limit of
TDDFT would not recover normal all-electron time-independent DFT.

These problems in handling nuclei plus electrons are evaded in the Li-Tong
\cite{LiTong86} multi-component version of RG TDDFT. By construction, of
course, the Li-Tong scheme has the same dependence on convergence of temporal
Taylor series expansions as the original RG scheme, hence cannot include
adiabatically switched external potentials. No full DFT dynamics (i.e.,
non-Kohn-Sham, even if purely formal) for the Li-Tong scheme has, to our
knowledge, ever been given. Though in principle general, the nuclear motion in
the Li-Tong approach in practice is analyzed purely at the KS level. See for
example Gross, Dobson, and Petersilka \cite{Gross96}. In the classical limit,
they obtain the correspondence between the KS dynamics for the model system
and the physical dynamics via the Ehrenfest theorem, with forces on the nuclei
from gradients of the TDKS potentials. Unexplored, as far as we are aware, is
the correspondence for either narrow wavepackets or full-quantum
nuclei.\cite{Xchneglect}

\section{Partitioned Equation of Motion in the $\mathbf{c}$ Parameterization}

Eq. (\ref{cgammaeom}), when rewritten in a form analogous with eq.
(\ref{dEdlambda}), has the obvious matrix partitioning
\begin{equation}
\left[
\begin{array}
[c]{cc}%
\eta_{\mathbf{cc}}\left(  t\right)  & \eta_{\mathbf{c}\gamma}\left(  t\right)
\\
\eta_{\gamma\mathbf{c}}\left(  t\right)  & \eta_{\gamma\gamma}\left(
t\right)
\end{array}
\right]  \left[
\begin{array}
[c]{c}%
\mathbf{\dot{c}}\\
\dot{\gamma}%
\end{array}
\right]  =\left[
\begin{array}
[c]{c}%
\nabla_{\mathbf{c}}\mathcal{E}\left(  t\right) \\
\nabla_{\gamma}\mathcal{E}\left(  t\right)
\end{array}
\right]  \label{partmat1}%
\end{equation}
This is an operator notation and $\nabla_{\lambda}\mathcal{E}\left(  t\right)
=\delta\mathcal{E}/\delta\lambda$. With $\lambda_{i}=\mathbf{c}\,${\ or
}$\gamma$, the components of the metric matrix $\eta$ are
\begin{equation}
\eta_{\lambda_{1},\lambda_{2}}(t,\,\mathbf{y},\,\mathbf{y}^{\prime})=i\left\{
\frac{\delta}{\delta\,\tilde{\lambda}_{1}(\mathbf{y})}\frac{\delta}%
{\delta\,\lambda_{2}(\mathbf{y}^{\prime})}\,-\,\frac{\delta}{\delta
\,\lambda_{2}(\mathbf{y})}\frac{\delta}{\delta\,\tilde{\lambda}_{1}%
(\mathbf{y}^{\prime})}\right\}  \ln S|_{\tilde{\lambda}_{i}=\lambda_{i}}
\label{etamat}%
\end{equation}
Here
\begin{equation}
S\equiv\left\langle \Psi(\tilde{\mathbf{c}}(t),\tilde{\gamma}(t))|\Psi
(\mathbf{c}(t),\gamma(t))\right\rangle
\end{equation}
and the analogue to the classical Hamiltonian (recall eq. (\ref{Eoft})) is
\begin{equation}
\mathcal{E}(t)=\frac{\left\langle \Psi(\mathbf{c}(t),\gamma(t))|H|\Psi
(\mathbf{c}(t),\gamma(t))\right\rangle }{\left\langle \Psi(\mathbf{c}%
(t),\gamma(t))|\Psi(\mathbf{c}(t),\gamma(t))\right\rangle } \label{Eofcgammat}%
\end{equation}
Coupled equations of motion follow
\begin{align}
\dot{\mathbf{c}}(t)  &  =\eta_{\mathbf{c\,c}}^{-1}(t)\left[  \nabla
_{\mathbf{c}}\mathcal{E}(t)-\eta_{\mathbf{c}\,\gamma}(t)\dot{\gamma}\right]
\nonumber\\
\dot{\gamma}(t)  &  =\left[  \eta_{\gamma\,\gamma}(t)-\eta_{\gamma
\,\mathbf{c}}(t)\eta_{\mathbf{c}\,\mathbf{c}}^{-1}(t)\eta_{\mathbf{c}\,\gamma
}(t)\right]  ^{-1}\nonumber\\
&  \times\left[  \nabla_{\gamma}\mathcal{E}(t)-\eta_{\gamma\,\mathbf{c}%
}(t)\eta_{\mathbf{c}\,\mathbf{c}}^{-1}(t)\nabla_{\mathbf{c}}\mathcal{E}%
(t)\right]  \label{cdotgammadot}%
\end{align}
Under the assumption that $\eta_{\mathbf{cc}}^{-1}(t)$ exists, denote by
$\mathbf{c}_{\#}(t)$ the solution of the first of these for the correct
density $\gamma_{\#}(t)$. Then the exact equation of motion for the temporal
development of the density is
\begin{equation}
\dot{\gamma}(t)=\tilde{\mathcal{N}_{1}}(\gamma(t))\left[  \nabla_{\gamma
}\,\mathcal{E}(\gamma(t))-\tilde{\mathcal{N}_{2}}(\gamma(t))\right]
\label{Appgammadot}%
\end{equation}
The functionals are defined as
%\begin{align}%
\begin{align}
\tilde{\mathcal{N}_{1}}(\gamma(t))  &  =\left[  \eta_{\gamma\,\gamma
}(\mathbf{c}_{\#}(t)\,,\gamma(t))-\eta_{\gamma\,\mathbf{c}}(\mathbf{c}%
_{\#}(t)\,,\gamma(t))\eta_{\mathbf{c}\,\mathbf{c}}^{-1}(\mathbf{c}%
_{\#}(t)\,,\gamma(t))\right.  \times\nonumber\\
&  \qquad\qquad\qquad\qquad\qquad\qquad\qquad\qquad\qquad\qquad\left.
\eta_{\mathbf{c}\,\gamma}(\mathbf{c}_{\#}(t)\,,\gamma(t))\right]
^{-1}\nonumber\\
\tilde{\mathcal{N}_{2}}(\gamma(t))  &  =\eta_{\gamma\,\mathbf{c}}%
(\mathbf{c}_{\#}(t)\,,\gamma(t))\eta_{\mathbf{c}\,\mathbf{c}}^{-1}%
(\mathbf{c}_{\#}(t)\,,\gamma(t))\nabla_{\mathbf{c}}\mathcal{E}(\mathbf{c}%
_{\#}(t)\,,\gamma(t)).\quad
\end{align}
Any singularities are associated with bifurcation points of the time evolution
of the system.

\section{Unitary Groups and Invariances of the Kernel}

\noindent Here we give details as to why it is necessary to use the general
linear group rather than the unitary group to describe the invariances of the
kernel. The time-dependence is inconsequential to the discussion, so is suppressed.

If consideration were restricted to unitary operators, the invariance group of
the kernel would be
\begin{equation}
\mathcal{G}_{\ker}^{U}=\left\{  U\in\mathcal{U}(\mathcal{H}^{N});\mathrm{Tr}%
\{\Phi^{\dagger}\Phi UYU^{\dagger}\}=\mathrm{Tr}\{\Phi^{\dagger}\Phi
Y\}=0\,\forall Y\in\ker\Xi\right\}  \label{Gkerdefn}%
\end{equation}
where $\mathcal{U}(\mathcal{H}^{N})$ is the group of unitary operators on
$N$-particle space. The condition expressed in eq. $\left(  \ref{Gkerdefn}%
\right)  $ does not imply that transformations belonging to $\mathcal{G}%
_{\ker}^{U}$ send elements of an equivalence class to itself
\begin{equation}
\mathcal{G}_{\ker}^{U}\left[  f\right]  _{N}\rightarrow\left[  f^{\prime
}\right]  _{N};\quad\text{in general }f\neq f^{\prime}%
\end{equation}
On the other hand the condition that defines the unitary group
\begin{equation}
\mathcal{G}_{\Xi}\ker\Xi\rightarrow\ker\Xi
\end{equation}
that does send equivalence classes to equivalence classes is given by
\begin{equation}
\mathcal{G}_{\Xi}=\left\{  U\in\mathcal{U}(\mathcal{H}^{N});\mathrm{Tr}%
\{\Phi^{\dagger}\Phi UYU^{\dagger}\}=\mathrm{Tr}\{\Phi^{\dagger}\Phi
Y\}\quad\,\forall Y\in\mathcal{B}_{1}\left(  \mathcal{H}^{N}\right)  \right\}
\label{GXi}%
\end{equation}
Evidently from eq. $\left(  \ref{GXi}\right)  $
\begin{equation}
\left[  U(\omega),\Phi^{\dagger}\Phi\right]  =0
\end{equation}
and $\mathcal{G}_{\Xi}$ is obviously a subgroup of $\mathcal{G}_{\ker}^{U}$
and can be represented by transformations of the form
\begin{align}
\mathcal{G}_{\Xi}  &  \subset\mathcal{G}_{\ker}^{U}=\{U(\omega);\hat{U}%
(\omega)(Y)\equiv U(\omega)YU^{\dagger}(\omega);\nonumber\\
U(\omega)  &  =e^{i\,\int\omega(\kappa_{1},\ldots,\kappa_{N})\Phi^{\dagger
}(\kappa_{1})\Phi(\kappa_{1})\ldots\Phi^{\dagger}(\kappa_{N})\Phi(\kappa
_{N})d\kappa_{1}\ldots d\kappa_{N}};\nonumber\\
&  \qquad\qquad\qquad\qquad\quad\omega(\kappa_{1},\ldots,\kappa_{N}%
)\in\mathbb{R}\} \label{GXidefn}%
\end{align}
One can see that $\mathcal{G}_{\Xi}$ is the largest commutative unitary
subgroup of $\mathcal{G}_{\ker}^{U}$ with these properties. Because of eq.
(\ref{WifromYref}), it follows that fibers associated with states are mapped
into themselves by $\hat{U}(\omega)$
\begin{equation}
Y\in\lbrack\gamma]_{N}\Leftrightarrow\hat{U}(\omega)Y\in\lbrack\gamma]_{N}%
\end{equation}
However, given a reference element $Y_{ref}\in\lbrack\gamma]_{N}$
corresponding to a pure state $|\Psi_{ref}\rangle$ in the fiber $[\gamma]_{N}%
$, it is \emph{not} the case that all the rest of the pure-state segment of
the fiber is generated by application of elements of $\mathcal{G}_{\Xi}$, (see
Appendix G), i.e. this group does \emph{not }act homogeneously on the set of
pure states in $[\gamma]_{N}$. This also shows that\emph{\ there is no intra
fiber unitary subgroup of }$\mathcal{G}_{\ker}^{U}$ \emph{that acts
homogeneously on the set of pure states contained within a fiber}.\emph{\ }The
physical insight is that a transformation such as eq. (\ref{GXidefn}) never
can generate the equivalent of a Jastrow or Hylleras factor in the Hilbert
space of pure states. If one drops the unitary requirement in eq. $\left(
\ref{Gkerdefn}\right)  $ \textit{it is} possible find transformations, that
are diagonal in the coordinate representation (like those of $\mathcal{G}%
_{\Xi}$)\textit{, }which produce all pure states with a given density,
$\gamma$, from a pure reference state $\left|  \Psi\right\rangle $ of that
density\textit{. }These transformations are of the form%
\begin{equation}
T\left(  \mathcal{V},\omega\right)  =e^{\int\left\{  \mathcal{V}(\kappa
_{1},\ldots,\kappa_{N})+i\omega(\kappa_{1},\ldots,\kappa_{N})\right\}
\Phi^{\dagger}(\kappa_{1})\Phi(\kappa_{1})\ldots\Phi^{\dagger}(\kappa_{N}%
)\Phi(\kappa_{N})d\kappa_{1}\ldots d\kappa_{N}}%
\end{equation}
where there are certain restriction placed on $\mathcal{V}(\kappa_{1}%
,\ldots,\kappa_{N})$ (see Appendix G). We denote the group formed by these
intra fiber transformations $\mathfrak{G}_{fib}$ in the main body of this
article. With a weaker set of restrictions, (see Appendix G), such maps also
transform $\ker\Xi\rightarrow\ker\Xi$, and form a subgroup $\mathfrak{G}%
_{\ker}\subset\mathfrak{G}_{\Xi}$, and with a third set of restrictions, one
forms a group of inter fiber transformations $\mathfrak{G}_{den}$ (see
Appendix F).

\section{Positive Maps of Trace Class Operators}

\noindent Recall that a trace class operator $Y \; \in\mathcal{B}_{1} (
\mathcal{H})$, with $\mathcal{H}$ any Hilbert space, is positive $Y \ge0$ if
\begin{equation}
\langle v | Y v \rangle\ge0 \qquad\forall\, v \in\mathcal{H}%
\end{equation}
which is true if and only if all the eigenvalues of $Y \ge0$. In turn this
requirement is true if and only if
\begin{equation}
Y = Q Q^{\dagger} \, , Q\in\mathcal{B}_{2} ( \mathcal{H})
\end{equation}
and $\mathcal{B}_{2} ( \mathcal{H})$ is the space of Hilbert-Schmidt operators.

A linear map of trace class operators $\hat{T}\mathcal{B}_{1}(\mathcal{H}%
)\mapsto\mathcal{B}_{1}(\mathcal{H})$ is positive if $\hat{T}(Y)\geq0$ for
$Y\geq0$. In particular, $\hat{T}(Y)=TYT^{\dagger}$ is a positive map
$\mathcal{B}_{1}(\mathcal{H})\mapsto\mathcal{B}_{1}(\mathcal{H})$ for any
bounded operator $T$, i.e., for such a $T$
\begin{equation}
\hat{T}(Y)\in\mathcal{B}_{1}(\mathcal{H})_{+}\,\,\text{{\ if }}\,Y\in
\mathcal{B}_{1}(\mathcal{H})_{+}%
\end{equation}
where $\mathcal{B}_{1}(\mathcal{H})_{+}$ is the cone of positive trace class
operators \cite{Kummer}. Cones facilitate analysis because they define order
relationships. That is, $A>B$ if $A-B\in\mathcal{B}_{1}(\mathcal{H})_{+}$. The
set of positive operators is a cone, not just a convex set, as normalization
does not have to be constant.

If $T$ is an invertible bounded operator, i.e. $T\in GL(\mathcal{H})$, then
\begin{equation}
T=e^{A}\,e^{iB}%
\end{equation}
with $A^{\dagger}=A$ and $B^{\dagger}=B$. Clearly $e^{iB}$ is unitary while
$e^{A}$ is invertible and positive. The set $\{e^{iB}\}$ forms the unitary
subgroup $\mathcal{U}(\mathcal{H})$ of $GL(\mathcal{H})$, while the set $\{
e^{A}\}$ comprises the left coset $GL(\mathcal{H})/\mathcal{U}(\mathcal{H})$
of $\mathcal{U}(\mathcal{H})$ in $GL(\mathcal{H})$. Denote the maps $\hat{T}$
generated by $T\in GL(\mathcal{H})$ as $\widehat{GL(\mathcal{H})}$ and observe
that they map $\mathcal{B}_{1}(\mathcal{H})\mapsto\mathcal{B}_{1}%
(\mathcal{H})$ and $\mathcal{B}_{2}(\mathcal{H})\mapsto\mathcal{B}%
_{2}(\mathcal{H}).$ These are used in the next Appendix.

\section{Group Orbits of Density Operators}

\noindent Once again the time-dependence appears only as a parameterization
that is inconsequential to the discussion, so is omitted. The transformation
\begin{align}
\hat{T}  &  \mathcal{B}_{1}(\mathcal{H})\mapsto\mathcal{B}_{1}\left(
\mathcal{H}\right) \nonumber\\
&  \hat{T}(Y)=TYT^{\dagger}%
\end{align}
maintains the rank of the operator $Y$, that is $\hat{T}(Y)$ has the same
number of non-zero eigenvalues as $Y$ itself. Particularizing to density
operators $D$, their cone $\mathcal{B}_{1}(\mathcal{H})_{+}$ can be classified
by the rank of $D\in\mathcal{B}_{1}(\mathcal{H})_{+}$. Operators for pure
states have unit rank and correspond to vectors in $\mathcal{H}$
\begin{equation}
\mathop{\rm rank}\nolimits(D)=1\,\,\Leftrightarrow D=\frac{|\Psi\rangle
\langle\Psi|}{\langle\Psi|\Psi\rangle}\,;|\Psi\rangle\in\mathcal{H}%
\end{equation}


All rank 1 density operators lie in the same orbit of $\widehat{GL(\mathcal{H}%
)} $ (see Appendix D for definition) in $\mathcal{B}_{1}(\mathcal{H})$. That
is, if $D_{0}$ is an arbitrary, specified pure state density operator, then
all other pure state density operators $D$ are generated by
\begin{equation}
D=\hat{T}(D_{0});\quad\hat{T}\in\widehat{GL(\mathcal{H})}\,;\quad
\mathop{\rm rank}\nolimits(D_{0})=1
\end{equation}


Ensemble states correspond to density operators with rank greater than unity.
As in the pure state case, ensemble density operators of a given rank all lie
on one orbit of $\widehat{GL(\mathcal{H})}$. Thus, if $D_{p,0}$ is an
arbitrary, specified ensemble state density operator of rank $p$, then all
other ensemble state density operators of the same rank $D_{p}$ are generated
by
\begin{equation}
D_{p}=\hat{T}(D_{p,0});\quad\hat{T}\in\widehat{GL(\mathcal{H})}\,;\quad
\mathop{\rm rank}\nolimits(D_{p,0})=p
\end{equation}


\section{Analytical Structure of the Kernel}

\noindent This Appendix treats characteristics of $Y\in\ker\Xi$. As before,
the time-dependence appears only as a parameterization that is inconsequential
to the discussion, so is omitted. It is useful for present purposes to begin
by expressing all operators belonging to $\mathcal{B}_{1}(\mathcal{H}^{N})$
(recall $N=N_{e}+N_{n}$) in terms of the ordered generalized operator basis
\cite{BerezinShubin}
\begin{equation}
\left\{  |\mathbf{x}_{1}\ldots\mathbf{x}_{N_{e}}\rangle_{ord}\langle
\mathbf{x}_{1}^{\prime}\ldots\mathbf{x}_{N_{e}}^{\prime}|_{ord}\otimes
|\mathbf{X}_{1}\ldots\mathbf{X}_{N_{n}}\rangle_{ord}\langle\mathbf{X}%
_{1}^{\prime}\ldots\mathbf{X}_{N_{n}}^{\prime}|_{ord}\right\}
\end{equation}
with
\begin{equation}
\mathbf{x},\mathbf{x}^{\prime},\mathbf{X},\mathbf{X}^{\prime}\in\mathbb{F}%
\end{equation}
where $\mathbb{F}=\mathbb{R}^{3}\otimes\mathbb{C}^{2}$ if the space-spin
analysis is pursued and $\mathbb{F}=\mathbb{R}^{3}$ if spin is suppressed
(whence the states are coordinate eigenfunctions). The subscript ``$ord$''
denotes the ordered basis
\begin{equation}
\mathbf{x}_{1}\leq\mathbf{x}_{2}\leq\mathbf{x}_{3}\leq\ldots
\end{equation}
with the understanding that strict inequality holds only for each species of
fermions. Shortly a shift to the unordered basis will be advantageous. For
either basis, a compressed notation, analogous with eq.~(\ref{compressed}), is
convenient
\begin{align}
&  |\mathbf{x}_{1}\ldots\mathbf{x}_{N_{e}}\rangle\langle\mathbf{x}_{1}%
^{\prime}\ldots\mathbf{x}_{N_{e}}^{\prime}|\otimes|\mathbf{X}_{1}%
\ldots\mathbf{X}_{N_{n}}\rangle\langle\mathbf{X}_{1}^{\prime}\ldots
\mathbf{X}_{N_{n}}^{\prime}|\nonumber\\
&  \equiv|\kappa_{1}\ldots\kappa_{N_{e}},\kappa_{N_{e}+1}\ldots\kappa
_{N}\rangle\langle\kappa_{1}^{\prime}\ldots\kappa_{N_{e}}^{\prime}%
,\kappa_{N_{e}+1}^{\prime}\ldots\kappa_{N}^{\prime}|\nonumber\\
&  \equiv|\mathsf{x},\mathsf{X}\rangle\langle\mathsf{x}^{\prime}%
,\mathsf{X}^{\prime}|,
\end{align}
The last form (recall eq. (\ref{manifdef})) will be used only when there is no ambiguity.

Any operator $Y\,\in\mathcal{B}_{1}(\mathcal{H}^{N})$ can be expanded in terms
of the ordered basis as
\begin{align}
Y  &  =\int\,d\kappa_{1},\ldots,d\kappa_{N},d\kappa_{1}^{\prime}%
,\ldots,d\kappa_{N}^{\prime}\tilde{Y}(\kappa_{1},\ldots,\kappa_{N};\kappa
_{1}^{\prime},\ldots,\kappa_{N}^{\prime})\nonumber\\
&  \times|\kappa_{1}\ldots\kappa_{N}\rangle_{ord}\langle\kappa_{1}^{\prime
}\ldots\kappa_{N}^{\prime}|_{ord}%
\end{align}
The operator integral kernel with respect to the ordered basis can be expanded
\cite{Richtmyer} as
\begin{align}
\tilde{Y}(\kappa_{1},\ldots,\kappa_{N},\kappa_{1}^{\prime},\ldots,\kappa
_{N}^{\prime})  &  =\sum_{i=1}^{\infty}\alpha_{i}\bar{\varphi}_{i}(\kappa
_{1},\ldots,\kappa_{N})\psi_{i}(\kappa_{1}^{\prime},\ldots,\kappa_{N}^{\prime
});\nonumber\\
\alpha_{i}  &  \geq0; \qquad\sum_{i=1}^{\infty}\alpha_{i}<\infty
\end{align}
where the basis functions $\left\{  \varphi_{i},\psi_{i}\right\}  $ are
specific to the operator $Y$ in the sense of spanning its range and domain and
satisfying
\begin{align}
\int\,d\kappa_{1},\ldots,d\kappa_{N},\bar{\varphi}_{i}(\kappa_{1}%
,\ldots,\kappa_{N})\varphi_{j}(\kappa_{1},\ldots,\kappa_{N})  &  =\delta
_{ij}\nonumber\\
\int\,d\kappa_{1},\ldots,d\kappa_{N},\bar{\psi}_{i}(\kappa_{1},\ldots
,\kappa_{N})\psi_{j}(\kappa_{1},\ldots,\kappa_{N})  &  =\delta_{ij}%
\end{align}


It is easier to consider the contraction map, the central quantity, in the
\emph{unordered} basis. In it, the operator integral kernel is
\begin{equation}
Y(\kappa_{1},\ldots,\kappa_{N},\kappa_{1}^{\prime},\ldots,\kappa_{N}^{\prime
})=\frac{(-1)^{(P_{e}+P_{n})}}{(N_{e}!N_{n}!)^{2}}\mathcal{P}\tilde{Y}%
(\kappa_{1},\ldots,\kappa_{N},\kappa_{1}^{\prime},\ldots,\kappa_{N}^{\prime})
\end{equation}
where $\mathcal{P}^{-1}$ is the permutation operator that generates the
ordered strings from the unordered ones
\begin{align}
\mathcal{P}^{-1}\{\kappa_{1},\ldots,\kappa_{N}\}  &  =\{\kappa_{1}%
,\ldots,\kappa_{N}\}_{ord}\nonumber\\
\mathcal{P}^{-1}\{\kappa_{1}^{\prime},\ldots,\kappa_{N}^{\prime}\}  &
=\{\kappa_{1}^{\prime},\ldots,\kappa_{N}^{\prime}\}_{ord}%
\end{align}
$P_{e}$ is the number of permutations of the electron coordinates, and $P_{n}
$ is either the number of permutations of the nuclear coordinates if the
nuclei are fermions or $0$ if they are bosons.

In the unordered basis the contraction map $\Xi$ eq. (\ref{Xidefn}) on any
operator $Y$ is
%\begin{align}%
\begin{align}
&  \!\!\!\!\!Tr\{\Phi^{\dagger}(\mathbf{x})\Phi(\mathbf{x})\Phi^{\dagger
}(\mathbf{X})\Phi(\mathbf{X})Y\}\nonumber\\
&  =\int\,d\kappa_{1},\ldots,d\kappa_{N}d\kappa_{1}^{\prime},\ldots
,d\kappa_{N}^{\prime}d\kappa_{1}^{\prime\prime},\ldots,d\kappa_{N}%
^{\prime\prime}\nonumber\\
&  \qquad\times\langle\kappa_{1}^{\prime\prime},\ldots,\kappa_{N}%
^{\prime\prime}|\Phi^{\dagger}(\mathbf{x})\Phi(\mathbf{x})\Phi^{\dagger
}(\mathbf{X})\Phi(\mathbf{X})|\kappa_{1},\ldots,\kappa_{N}\rangle\nonumber\\
&  \qquad\times\langle\kappa_{1}^{\prime},\ldots,\kappa_{N}^{\prime}%
|\kappa_{1}^{\prime\prime},\ldots,\kappa_{N}^{\prime\prime}\rangle
\;Y(\kappa_{1},\ldots,\kappa_{N};\kappa_{1}^{\prime},\ldots,\kappa_{N}%
^{\prime})\nonumber\\
&  =N_{e}N_{n}\,\int d\mathbf{x}_{2}^{\prime}\ldots,d\mathbf{x}_{N_{e}%
}^{\prime}d\mathbf{X}_{2}^{\prime},\ldots,d\mathbf{X}_{N_{n}}^{\prime
}\nonumber\\
&  \times Y(\mathbf{x},\mathbf{x}_{2}\ldots,\mathbf{x}_{N_{e}}\mathbf{X}%
,\mathbf{X}_{2}\ldots,\mathbf{X}_{N_{n}};\mathbf{x},\mathbf{x}_{2}^{\prime
}\ldots,\mathbf{x}_{N_{e}}^{\prime}\mathbf{X},\mathbf{X}_{2}^{\prime}%
\ldots,\mathbf{X}_{N_{n}}^{\prime}) \label{TrphiphiphiphiY}%
\end{align}


The elements of $\ker\Xi$ become evident by writing
%\begin{align}%
\begin{align}
Y(\kappa_{1},\ldots,\kappa_{N};\kappa_{1}^{\prime},\ldots,\kappa_{N}^{\prime
})  &  =Y_{d}(\kappa_{1},\ldots,\kappa_{N})\prod_{j=1}^{N}\delta(\kappa
_{j}-\kappa_{j}^{\prime})\nonumber\\
&  +Y_{OD}(\kappa_{1},\ldots,\kappa_{N};\kappa_{1}^{\prime},\ldots,\kappa
_{N}^{\prime}) \label{YYdYod}%
\end{align}
\begin{align}
Y_{OD}(\kappa_{1},\ldots,\kappa_{N};\kappa_{1},\ldots,\kappa_{N})  &
=0\nonumber\\
Y_{d}(\kappa_{1},\ldots,\kappa_{N})  &  \equiv Y(\kappa_{1},\ldots,\kappa
_{N};\kappa_{1},\ldots,\kappa_{N})
\end{align}
and the caveat that the product of delta functions in eq. (\ref{YYdYod})
implicitly includes all the permutations of coordinates (because of working in
the unordered basis). Note that $Y_{d}$ is symmetric with respect to
interchange among the electron coordinates themselves and similarly for the
nuclear coordinates. Substitution of (\ref{YYdYod}) in (\ref{TrphiphiphiphiY})
gives
%\begin{gather}
%Tr\!\{\Phi ^{\dagger }({\bf x})\Phi ({\bf x})\Phi ^{\dagger }({\bf X})\Phi (%
%{\bf X})Y\}=  \nonumber \\
%N_{e}N_{n}\,\int d{\bf x}_{2}^{\prime }\ldots ,d{\bf x}_{N_{e}}^{\prime }d%
%{\bf X}_{2}^{\prime },\ldots ,d{\bf X}_{N_{n}}^{\prime }Y_{d}({\bf x},{\bf x}%
%_{2}\ldots ,{\bf x}_{N_{e}}{\bf X},{\bf X}_{2}\ldots ,{\bf X}_{N_{n}})
%\label{ContrYtoYd}
%\end{gather}%
An alternative way to state this is that the contraction map sends all
off-diagonal basis operators to zero
\begin{align}
&  \Xi\left\{  |\kappa_{1}\ldots\kappa_{N}\rangle\langle\kappa_{1}^{\prime
}\ldots\kappa_{N}^{\prime}|\right\}  =0\nonumber\\
&  \;\text{{for }}\left\{  \kappa_{1}\ldots\kappa_{N_{e}}\right\}
\neq\left\{  \kappa_{1}^{\prime}\ldots\kappa_{N_{e}}^{\prime}\right\}  \text{
or }\left\{  \kappa_{1}\ldots\kappa_{N_{n}}\right\}  \neq\left\{  \kappa
_{1}^{\prime}\ldots\kappa_{N_{n}}^{\prime}\right\}  \label{zeroeta}%
\end{align}


Three cases are obvious immediately from eq. (\ref{ContrYtoYd}); two belong to
$\ker\Xi$, one does not. First, there are pure off-diagonal operators
\begin{equation}
Y=Y_{OD}\,\,;\,\,(Y_{d}=0)\,\Rightarrow Y\in\ker\Xi\label{Ycase1}%
\end{equation}
Then there are diagonal operators with a special property
%\begin{gather}
%Y=Y_{d};\,\,(Y_{OD}=0)  \nonumber \\
%\int d{\bf x}_{2}\ldots ,d{\bf x}_{N_{e}}d{\bf X}_{2},\ldots ,d{\bf X}%
%_{N_{n}}Y_{d}({\bf x},{\bf x}_{2}\ldots ,{\bf x}_{N_{e}}{\bf X},{\bf X}%
%_{2}\ldots ,{\bf X}_{N_{n}})=0  \nonumber \\
%\Rightarrow Y\in \ker \Xi  \label{Ycase2}
%\end{gather}%
Finally, there are diagonal operators that do not belong to $\ker\Xi$
%\begin{gather}
%Y=Y_{d};\,\,(Y_{OD}=0)  \nonumber \\
%0\neq \int d{\bf x}_{2}\ldots ,d{\bf x}_{N_{e}}d{\bf X}_{2},\ldots ,d{\bf X}%
%_{N_{n}}Y_{d}({\bf x},{\bf x}_{2}\ldots ,{\bf x}_{N_{e}}{\bf X},{\bf X}%
%_{2}\ldots ,{\bf X}_{N_{n}})  \nonumber \\
%\Rightarrow Y\notin \ker \Xi  \label{Ycase3}
%\end{gather}


The space $K_{OD}$ generated by operators satisfying (\ref{Ycase1}) is a
subspace of $\ker\{\Xi\}$
\begin{align}
K_{OD}  &  \subset\ker\{\Xi\}\nonumber\\
K_{OD}  &  =\mathrm{linear\,span}\left\{  Y=Y_{OD},Y_{d}=0\right\}
\label{koddefn}%
\end{align}
The remaining part $K_{D}$ of $\ker\{\Xi\}$ is comprised of diagonal operators
that satisfy eq. (\ref{Ycase2})
\begin{align}
K_{D}  &  \subset\ker\{\Xi\};\text{ where }K_{D}=\mathrm{linear\,span}\left\{
Y=Y_{d},Y_{OD}=0;\right. \nonumber\\
&  \left.  0=\int d\mathbf{x}_{2},\ldots d\mathbf{x}_{N_{e}}d\mathbf{X}%
_{2},\ldots d\mathbf{X}_{N_{n}}Y_{d}(\mathbf{x},\mathbf{x}_{2},\ldots
\mathbf{x}_{N_{e}}\mathbf{X},\mathbf{X}_{2},\ldots\mathbf{X}_{N_{n}})\right\}
\nonumber\\
&  \label{kddefn}%
\end{align}
In detail, for $K_{D}$, $Y_{d}$ obeys at least one of the two conditions
\begin{align}
0  &  =\int\,d\mathbf{x}_{i_{1}},\ldots,d\mathbf{x}_{i_{N_{e}-1}}%
,Y_{d}(\mathbf{x}_{1},\ldots,\mathbf{x}_{N_{e}}\mathbf{X}_{1},\ldots
,\mathbf{X}_{N_{n}})\nonumber\\
0  &  =\int\,d\mathbf{X}_{j_{1}},\ldots,d\mathbf{X}_{j_{N_{n}-1}}%
,Y_{d}(\mathbf{x}_{1},\ldots,\mathbf{x}_{N_{e}}\mathbf{X}_{1},\ldots
,\mathbf{X}_{N_{n}}) \label{kddefn1}%
\end{align}
with
\begin{align}
\left\{  i_{1}\ldots i_{N_{e}-1}\right\}   &  \subset\left\{  1\ldots
N_{e}\right\} \nonumber\\
\left\{  j_{1}\ldots j_{N_{n}-1}\right\}   &  \subset\left\{  1\ldots
N_{n}\right\}  \label{indices}%
\end{align}


Thus the kernel of $\Xi$ can be expressed as
\begin{equation}
\ker\{\Xi\}=K_{OD}\oplus K_{D}%
\end{equation}
with the two components orthogonal with respect to the Hilbert-Schmidt inner
product
\begin{equation}
Tr\{k_{OD}^{\dagger}k_{D}\}=0\;\mathrm{or}\,\,k_{OD}\in K_{OD},\,\,k_{d}\in
K_{D}%
\end{equation}
This result follows from
\begin{align}
Tr\{|\kappa_{1}\ldots\kappa_{N}\rangle\langle\kappa_{1}\ldots\kappa
_{N}||\kappa_{1}^{\prime}\ldots\kappa_{N}^{\prime}\rangle\langle\kappa
_{1}^{\prime\prime}\ldots\kappa_{N}^{\prime\prime}|\}  &  =\nonumber\\
\langle\kappa_{1}\ldots\kappa_{N}|\kappa_{1}^{\prime}\ldots\kappa_{N}^{\prime
}\rangle\langle\kappa_{1}^{\prime\prime}\ldots\kappa_{N}^{\prime\prime}%
|\kappa_{1}\ldots\kappa_{N}\rangle &  =0
\end{align}
which holds when both
\begin{equation}
\{\kappa_{1}^{\prime}\ldots\kappa_{N_{e}}^{\prime}\}\neq\{\kappa_{1}%
^{\prime\prime}\ldots\kappa_{N_{e}}^{\prime\prime}\}
\end{equation}
\emph{and}
\begin{equation}
\{\kappa_{N_{e}+1}^{\prime}\ldots\kappa_{N}^{\prime}\}\neq\{\kappa_{N_{e}%
+1}^{\prime\prime}\ldots\kappa_{N}^{\prime\prime}\}
\end{equation}
In either case the primed and double primed sets of coordinates have at least
one difference, hence both cannot be equal to $\{\kappa_{1}\ldots\kappa_{N}\}$.

Eq. (\ref{Ycase3}) makes clear that the only operators $Y\in\mathcal{B}%
_{1}(\mathcal{H}^{N})$ that do not lie in $\ker\Xi$ are diagonal operators
that do not belong to $K_{D}$. These form the subspace (of $\mathcal{B}%
_{1}(\mathcal{H}^{N})$)
\begin{align}
S_{D}  &  =\mathrm{linear\,span}\big\{Y=Y_{d},Y_{OD}=0;\big.\qquad\nonumber\\
&  \big.0\neq\!\!\!\int d\mathbf{x}_{2}\ldots,d\mathbf{x}_{N_{e}}%
d\mathbf{X}_{2},\ldots,d\mathbf{X}_{N_{n}}Y_{d}(\mathbf{x},\mathbf{x}%
_{2}\ldots,\mathbf{x}_{N_{e}}\mathbf{X},\mathbf{X}_{2}\ldots,\mathbf{X}%
_{N_{n}})\!\big\}\nonumber\\
&  \label{sddefn}%
\end{align}
Therefore the space of trace class operators has the direct sum decomposition
\begin{equation}
\mathcal{B}_{1}(\mathcal{H}^{N})=K_{OD}\oplus K_{D}\oplus S_{D} \label{decomp}%
\end{equation}


\section{Vector Bundle Transformations}

\noindent This Appendix develops pertinent properties of transformations
$TYT^{\dagger}$ of trace class operators $Y$ in $K_{OD}\oplus K_{D}$ $\oplus$
$S_{D}$ by elements $T$ of the General Linear Group $GL\left(  \mathcal{H}%
^{N}\right)  $ of $\mathcal{H}^{N}$. In the coordinate eigenfunction basis,
the transformation $TYT^{\dagger}$ reads
\begin{equation}
TYT^{\dagger}\equiv\hat{T}Y=\int d\mathsf{z}d\mathsf{z}^{\prime}\left[
\,TYT^{\dagger}(\mathsf{z},\mathsf{z}^{\prime})\,\right]  |\mathsf{z}%
\rangle\langle\mathsf{z}^{\prime}| \label{TYT}%
\end{equation}
where the transformed integral kernel is
\begin{equation}
TYT^{\dagger}(\mathsf{z},\mathsf{z}^{\prime})=\int d\mathsf{z}^{\prime\prime
}d\mathsf{z}^{\prime\prime\prime}T(\mathsf{z},\mathsf{z}^{\prime\prime
})Y(\mathsf{z}^{\prime\prime},\mathsf{z}^{\prime\prime\prime})T^{\ast
}(\mathsf{z}^{\prime\prime\prime},\mathsf{z}^{\prime}) \label{TYTkernel}%
\end{equation}
Clearly the only transforms that can preserve the structure described by eqs.
(\ref{koddefn}), (\ref{kddefn}), (\ref{sddefn}) have diagonal representations
\begin{align}
T(\mathsf{z},\mathsf{z}^{\prime})  &  =\delta(\mathsf{z}-\mathsf{z}^{\prime
})\,T(\mathsf{z})\nonumber\\
T(\mathsf{z})  &  =T(\mathsf{x},\mathsf{X}) \label{Tddefn}%
\end{align}
with $T(\mathsf{z})$ a species symmetric function with respect to interchange
among electronic coordinates and similarly for nuclear coordinates (within
each species).

The explicit form of the diagonal operator is
\begin{align}
T  &  =e^{\int\mathcal{V}(\mathbf{x}_{1},\ldots,\mathbf{x}_{N_{e}}%
\mathbf{X}_{1},\ldots,\mathbf{X}_{N_{n}})\prod\limits_{j=1}^{N_{e}}\Phi
_{e}^{\dagger}(x_{j})\Phi_{e}(x_{j})\prod\limits_{k=1}^{N_{n}}\Phi
_{n}^{\dagger}(X_{k})\Phi_{n}(X_{k})\prod\limits_{l=1}^{N_{e}}dx_{l}%
\prod_{m=1}^{N_{n}}dX_{m}}\nonumber\\
&  \;\times e^{i\int\omega(\mathbf{x}_{1},\ldots,\mathbf{x}_{N_{e}}%
\mathbf{X}_{1},\ldots,\mathbf{X}_{N_{n}})\prod\limits_{q=1}^{N_{e}}\Phi
_{e}^{\dagger}(x_{q})\Phi_{e}(x_{q})\prod\limits_{r=1}^{N_{n}}\Phi
_{n}^{\dagger}(X_{r})\Phi_{n}(X_{r})\prod\limits_{s=1}^{N_{e}}dx_{s}%
\prod_{t=1}^{N_{n}}dX_{t}} \label{Tdexplicit}%
\end{align}
When $T$ operates upon an element of the basis of many-body coordinate
eigenfunctions, the result is
\begin{align}
&  T\left|  y_{1}\ldots y_{N_{e}}Y_{1}\ldots Y_{N_{n}}\right\rangle
\nonumber\\
&  \quad=e^{\mathcal{V}(y_{1}\ldots y_{N_{e}}Y_{1}\ldots Y_{N_{n}})}%
e^{i\omega(y_{1}\ldots y_{N_{e}}Y_{1}\ldots Y_{N_{n}})}\left|  y_{1}\ldots
y_{N_{e}}Y_{1}\ldots Y_{N_{n}}\right\rangle \label{Tdcoord}%
\end{align}
because
\begin{align}
\Phi_{e}^{\dagger}(x_{j})\Phi_{e}(x_{j})\left|  y_{1}\ldots y_{N_{e}}%
Y_{1}\ldots Y_{N_{n}}\right\rangle  &  =\left|  y_{1}\ldots y_{N_{e}}%
Y_{1}\ldots Y_{N_{n}}\right\rangle \nonumber\\
\qquad\qquad\qquad\qquad\qquad\qquad\qquad\qquad\text{{for}}\;x_{j}  &
\in\{y_{1}\ldots y_{N_{e}}\}\nonumber\\
&  =0\,\,\text{o{therwise}}%
\end{align}
and the corresponding relationship for $\Phi_{n}^{\dagger}(X_{k})\Phi
_{n}(X_{k})$.

For $Y=Y_{K_{OD}}\,\in K_{OD}$, diagonal transformations $\hat{T}$ map
$K_{OD}\mapsto K_{OD}$ with the explicit form of the transformed integral
kernel
\begin{equation}
TY_{K_{OD}}T^{\dagger}(\mathsf{z},\mathsf{z}^{\prime})=T(\mathsf{z})Y_{K_{OD}%
}(\mathsf{z},\mathsf{z}^{\prime})T^{\ast}(\mathsf{z}^{\prime})\in K_{OD}
\label{TYodTkernel}%
\end{equation}
Diagonal transformations that map $K_{D}\mapsto K_{D}$ provided that the
transformed integral kernel satisfies eq. (\ref{Ycase2}), which is in fact the
defining property for maps that maintain the vector bundle structure of
$\mathfrak{B}_{vec}$. This can be written as
\begin{align}
\hat{T}  &  Y_{K_{D}}\mapsto TY_{K_{D}}T^{\dagger};\quad Y_{K_{D}}\in
K_{D}\nonumber\\
&  TY_{K_{D}}T^{\dagger}(\mathsf{z})=|T(\mathsf{z})|^{2}Y_{K_{D}}(\mathsf{z})
\end{align}
\emph{provided} that
\begin{equation}
|T(\mathsf{z})|^{2}=1+\phi(\mathsf{z}) \label{phidef}%
\end{equation}
and the species symmetric functions $\phi$ have the property
%\begin{align}%
\begin{align}
0  &  =\int\,d\mathbf{x}_{i_{1}},\ldots,d\mathbf{x}_{i_{N_{e}-1}}%
,d\mathbf{X}_{j_{1}},\ldots,d\mathbf{X}_{j_{N_{n}-1}}\nonumber\\
&  \times\phi(\mathbf{x}_{1},\ldots,\mathbf{x}_{N_{e}}\mathbf{X}_{1}%
,\ldots,\mathbf{X}_{N_{n}})Y_{d}(\mathbf{x}_{1},\ldots,\mathbf{x}_{N_{e}%
}\mathbf{X}_{1},\ldots,\mathbf{X}_{N_{n}})\nonumber\\
&  \qquad\qquad\qquad\qquad\qquad\qquad\qquad\qquad\qquad\quad\forall
Y_{K_{D}}\in K_{D} \label{strongorthog}%
\end{align}
and the indices on the variables of integration obey eq. (\ref{indices}).
These conditions follow directly from eqs. (\ref{Ycase2}). and
(\ref{strongorthog}) and expresses the property of strong orthogonality,
denoted by $\breve{\phi}\perp_{S}K_{D}$, or equivalently,
\begin{equation}
\breve{\phi}\perp_{S}\forall Y_{K_{D}}\in K_{D}%
\end{equation}
where $\breve{\phi}$ is the operator defined by the integral kernel $\phi$. In
particular, eq. (\ref{strongorthog}) implies that the pointwise product
function $\phi Y_{d}\in L_{1}(\mathtt{Z})$, which in turn implies that
$\phi\in L_{\infty}(\mathtt{Z})$ since $Y(\mathsf{z})\in L_{1}(\mathtt{Z})$.
Here
\begin{equation}
\mathtt{Z}=\mathbb{R}^{3N_{e}}\times\mathbb{R}^{3N_{n}}\times\mathbb{C}%
^{2N_{e}}\times\mathbb{C}^{2N_{n}}%
\end{equation}


Vector bundle transformations built from $\phi$ satisfy
%\begin{gather}
%\phi \in L_{\infty }({\tt Z})  \nonumber \\
%\phi \perp _{S}K_{D}  \label{phiprop}
%\end{gather}%
and have the property of mapping $K_{D}\mapsto K_{D}$, $S_{D}\mapsto
K_{D}\oplus S_{D}$ and form the group $\mathcal{G}_{\ker}$. This group
together with the unitary group $\mathcal{G}_{\Xi}$ forms the group
$\mathfrak{G}_{\ker}=\mathcal{G}_{\ker}\times\mathcal{G}_{\Xi}$, which is a
subgroup of the commutative subgroup $\mathcal{C}(\mathcal{H}^{N})\subset
GL(\mathcal{H}^{N})$. The elements of $\mathcal{C}(\mathcal{H}^{N})$ are given
by
\begin{equation}
T=e^{\left\{  \mathcal{\breve{V}}+i\breve{\omega}\right\}  }=T_{d}%
e^{i\breve{\omega}}%
\end{equation}
where the operator $T_{d}$ is defined by the integral kernel
\begin{equation}
T_{d}(\kappa_{1}\ldots\kappa_{N})=\exp[\mathcal{V}(\kappa_{1}\ldots\kappa
_{N})]=(1+\phi(\kappa_{1}\ldots\kappa_{N}))^{\frac{1}{2}}%
\end{equation}
with no restrictions on $\phi$. The integral kernel of $\mathcal{\breve{V}}$
is given by
\begin{equation}
\mathcal{V}(\kappa_{1}\ldots\kappa_{N})\equiv\frac{1}{2}\ln(1+\phi(\kappa
_{1}\ldots\kappa_{N})) \label{Vhalfln}%
\end{equation}
The transformations $\hat{T}K_{OD}\mapsto K_{OD}$ are determined by elements
of $\mathfrak{G}_{\ker}$ and act according to eq. $\left(  \ref{TYodTkernel}%
\right)  $.

We now develop the conditions $\phi$ must satisfy to determine the subgroups
$\mathfrak{G}_{den}$ and $\mathfrak{G}_{fib}$ of inter-fiber and intra-fiber
maps respectively contained in $\mathfrak{G}_{\ker}$. Once this is
accomplished we can construct a non-redundant parameterizations of pure states
in the equivalence classes $\left[  \gamma\right]  _{N}$ which can be utilized
in the constrained search procedure.

First consider transformations, $\hat{T}_{den}\left(  \phi\right)  $, between
fibers that produce the subgroup $\mathfrak{G}_{den}$, eqs. (\ref{phidef}),
(\ref{phiprop}) give the conditions for $\phi(\kappa_{1}\ldots\kappa_{N})$ to
map $S_{D}\mapsto S_{D}$. In that case
\begin{align}
&  \int\phi(\mathbf{x}_{1},\ldots,\mathbf{x}_{N_{e}}\mathbf{X}_{1}%
,\ldots,\mathbf{X}_{N_{n}})Y_{d}(\mathbf{x}_{1},\ldots,\mathbf{x}_{N_{e}%
}\mathbf{X}_{1},\ldots,\mathbf{X}_{N_{n}})\nonumber\\
&  \qquad\qquad\qquad\qquad\times d\mathbf{x}_{2},\ldots,d\mathbf{x}_{N_{e}%
},d\mathbf{X}_{2},\ldots,d\mathbf{X}_{N_{n}}\nonumber\\
&  =\int\phi(\mathbf{x}_{1},\ldots,\mathbf{x}_{N_{e}}\mathbf{X}_{1}%
,\ldots,\mathbf{X}_{N_{n}})Y_{S_{D}}(\mathbf{x}_{1},\ldots,\mathbf{x}_{N_{e}%
}\mathbf{X}_{1},\ldots,\mathbf{X}_{N_{n}})\nonumber\\
&  \qquad\qquad\qquad\qquad\times d\mathbf{x}_{2},\ldots,d\mathbf{x}_{N_{e}%
},d\mathbf{X}_{2},\ldots,d\mathbf{X}_{N_{n}}\neq0
\end{align}
with%
\begin{equation}
Y_{d}=Y_{K_{D}}+Y_{S_{D}} \label{YKdYSd}%
\end{equation}
This corresponds to
\begin{equation}
\Phi_{S_{D}}\left(  \mathsf{z}_{1},\mathsf{z}_{2}\right)  =\int P_{S_{D}%
}\left(  \mathsf{z}_{1},\mathsf{z}_{1}^{\prime}\right)  \phi\left(
\mathsf{z}_{1}^{\prime}\right)  P_{S_{D}}\left(  \mathsf{z}_{1}^{\prime
},\mathsf{z}_{2}\right)  d\mathsf{z}_{1}^{\prime}\neq0\quad\forall
\mathsf{z}_{1},\mathsf{z}_{2} \label{projS}%
\end{equation}
and hence
\begin{verbatim}
$\phi \NEG{\bot}_{s}Y_{d}$,
\end{verbatim}

where $P_{S_{D}}$ is the integral kernel of the projector onto the subspace
$S_{D}$. In order to obtain the defining representation of $\mathfrak{G}%
_{den}$ on $K_{D}\oplus S_{D}$, the functions $\phi$ should have zero effect
on $K_{D}$ thus need to satisfy
\begin{equation}
\Phi_{K_{D}}\left(  \mathsf{z}_{1},\mathsf{z}_{2}\right)  =\int P_{K_{D}%
}\left(  \mathsf{z}_{1},\mathsf{z}_{1}^{\prime}\right)  \phi\left(
\mathsf{z}_{1}^{\prime}\right)  P_{K_{D}}\left(  \mathsf{z}_{1}^{\prime
},\mathsf{z}_{2}\right)  d\mathsf{z}_{1}^{\prime}=0\quad\forall\mathsf{z}%
_{1},\mathsf{z}_{2} \label{projK}%
\end{equation}
and
\begin{equation}
\Phi_{K_{D}S_{D}}\left(  \mathsf{z}_{1},\mathsf{z}_{2}\right)  =\int P_{K_{D}%
}\left(  \mathsf{z}_{1},\mathsf{z}_{1}^{\prime}\right)  \phi\left(
\mathsf{z}_{1}^{\prime}\right)  P_{S_{D}}\left(  \mathsf{z}_{1}^{\prime
},\mathsf{z}_{2}\right)  d\mathsf{z}_{1}^{\prime}=0\quad\forall\mathsf{z}%
_{1},\mathsf{z}_{2} \label{projKS}%
\end{equation}
The effect of these transformations on the off diagonal part of $\mathcal{B}%
_{1}\left(  \mathcal{H}^{N}\right)  $, are given by elements of $\mathcal{G}%
_{\Xi}$ that act according to eq. $\left(  \ref{TYodTkernel}\right)  $ and is
given by
\begin{equation}
\hat{T}_{den}\left(  \phi\right)  Y_{K_{OD}}=e^{\mathcal{\breve{V}}%
_{den}\left(  \phi\right)  }Y_{K_{OD}}e^{\mathcal{\breve{V}}_{den}\left(
\phi\right)  }%
\end{equation}


In order to produce intra-fiber transformations, $\hat{T}_{fib}\left(
\phi\right)  K_{D}\mapsto K_{D}$, that form the subgroup $\mathcal{G}_{fib}$
the functions $\phi$ must satisfy the condition
\begin{equation}
\Phi_{S_{D}}\left(  \mathsf{z}_{1},\mathsf{z}_{2}\right)  =\int P_{S_{D}%
}\left(  \mathsf{z}_{1},\mathsf{z}_{1}^{\prime}\right)  \phi\left(
\mathsf{z}_{1}^{\prime}\right)  P_{S_{D}}\left(  \mathsf{z}_{1}^{\prime
},\mathsf{z}_{2}\right)  d\mathsf{z}_{1}^{\prime}=0\quad\forall\mathsf{z}%
_{1},\mathsf{z}_{2}%
\end{equation}
The defining representation of $\mathcal{G}_{fib}$ on $K_{D}\oplus S_{D}$, is
obtained when the functions further satisfy
\begin{equation}
\Phi_{K_{D}S_{D}}\left(  \mathsf{z}_{1},\mathsf{z}_{2}\right)  =\int P_{K_{D}%
}\left(  \mathsf{z}_{1},\mathsf{z}_{1}^{\prime}\right)  \phi\left(
\mathsf{z}_{1}^{\prime}\right)  P_{S_{D}}\left(  \mathsf{z}_{1}^{\prime
},\mathsf{z}_{2}\right)  d\mathsf{z}_{1}^{\prime}=0\quad\forall\mathsf{z}%
_{1},\mathsf{z}_{2}%
\end{equation}
The effect of these transformations on the off diagonal subspace $K_{OD}$ and
is given by
\begin{equation}
\hat{T}_{fib}\left(  \phi\right)  Y_{K_{OD}}=e^{\mathcal{\breve{V}}%
_{fib}\left(  \phi\right)  }Y_{K_{OD}}e^{\mathcal{\breve{V}}_{fib}\left(
\phi\right)  }%
\end{equation}
The group of intra fiber transformations that acts homogeneously on the set of
pure states within a fiber is $\mathfrak{G}_{fib}=\mathcal{G}_{fib}%
\times\mathcal{G}_{\Xi}$ i.e is constructed from transformations of the form
$e^{\mathcal{\breve{V}}_{fib}\left(  \phi\right)  }e^{i\breve{\omega}}$

\section{Non-Redundant Parameterization of Pure States Adapted to
$\mathfrak{B}_{vec}$}

In order to produce non-redundant parameterizations of pure states in
$\mathcal{B}_{1}\left(  \mathcal{H}^{N}\right)  $ by the groups $\mathfrak{G}%
_{fib}$ and $\mathfrak{G}_{den}$ one needs to introduce biorthogonal bases for
$K_{OD},K_{D}$ and $S_{D}$. (The space of trace class operators $\mathcal{B}%
_{1}\left(  \mathcal{H}^{N}\right)  $ is not a Hilbert space thus is not self
dual, its dual space is the space of bounded operators $\mathcal{B}\left(
\mathcal{H}^{N}\right)  )$.The basis for $K_{OD},K_{D}$ and $S_{D}$ are
\begin{align}
&  K_{OD}=\mathrm{linear\,span}\left\{  \breve{h}_{\iota}\breve{h}_{\iota
}=\int h_{\iota}\left(  \mathsf{z}\mathbf{,}\mathsf{z}^{\prime}\right)
\left|  \mathsf{z}\right\rangle \left\langle \mathsf{z}^{\prime}\right|
d\mathsf{z}d\mathsf{z}^{\prime};\mathsf{z}\mathbf{\neq}\mathsf{z}^{\prime
}\right\} \\
&  K_{D}=\mathrm{linear\,span}\left\{  \breve{k}_{\iota}\breve{k}_{\iota}=\int
k_{\iota}\left(  \mathsf{z}\right)  \left|  \mathsf{z}\right\rangle
\left\langle \mathsf{z}\right|  d\mathsf{z}\right. \nonumber\\
&  \left.  \int k_{\iota}\left(  \kappa_{2,}..,\kappa_{N_{e}},\kappa_{N_{e}%
+1},..\kappa_{N}\right)  d\kappa_{2}\ldots d\kappa_{N_{e}}d\kappa_{N_{e}%
+2}\ldots d\kappa_{N}\mathbf{=0}\text{ }\mathbf{a.e.}\right\} \label{kbasis}\\
&  S_{D}=\mathrm{linear\,span}\left\{  \breve{s}_{\iota}\breve{s}_{\iota}=\int
s_{\iota}\left(  \mathsf{z}\right)  \left|  \mathsf{z}\right\rangle
\left\langle \mathsf{z}\right|  d\mathsf{z}\right.  ;\nonumber\\
&  \left.  \int s_{\iota}\left(  \kappa_{2,}..,\kappa_{N_{e}},\kappa_{N_{e}%
+1},..\kappa_{N}\right)  d\kappa_{2}\ldots d\kappa_{N_{e}}d\kappa_{N_{e}%
+2}\ldots d\kappa_{N}\mathbf{\neq0\,a.e.}\right\}  \label{sbasis}%
\end{align}
where $\mathbf{a.e.\equiv}$ ''almost everywhere'' i.e. except on sets of
measure zero. These bases are biothorgonal to the bases $\left\{  \breve
{k}_{\kappa}^{d},\breve{s}_{\kappa}^{d};1\leq\kappa\leq\infty\right\}  $ of
$\mathcal{B}\left(  \mathcal{H}^{N}\right)  $ i.e
\begin{align}
\left(  \breve{h}_{\kappa}^{d}|\breve{h}_{\iota}\right)   &  =\left(
\breve{k}_{\kappa}^{d}|\breve{k}_{\iota}\right)  =\left(  \breve{s}_{\kappa
}^{d}|\breve{s}_{\iota}\right)  =\delta_{\kappa\iota}\nonumber\\
\left(  \breve{k}_{\kappa}^{d}|\breve{s}_{\iota}\right)   &  =\left(
\breve{s}_{\kappa}^{d}|\breve{k}_{\iota}\right)  =\left(  \breve{h}_{\kappa
}^{d}|\breve{s}_{\iota}\right) \nonumber\\
&  =\left(  \breve{h}_{\kappa}^{d}|\breve{k}_{\iota}\right)  =\left(
\breve{k}_{\kappa}^{d}|\breve{h}_{\iota}\right)  =\left(  \breve{s}_{\kappa
}^{d}|\breve{h}_{\iota}\right)  =0
\end{align}
and we have introduced the super operator scalar product $\left(  A|B\right)
=Tr\left\{  A^{\dagger}B\right\}  $.

If we are given a reference pure state $D_{ref}$ that has a density $\gamma$
then it can be decomposed as
\begin{equation}
D_{ref}^{N}{}_{K_{OD}}+D_{ref}^{N}{}_{K_{D}}+\Gamma\left(  \gamma\right)
\end{equation}
and we can define
\begin{align}
\breve{h}_{1}  &  =D_{ref}^{N}{}_{K_{OD}}\nonumber\\
\breve{k}_{1}  &  =D_{ref}^{N}{}_{K_{D}}\nonumber\\
\breve{s}_{1}  &  =\Gamma\left(  \gamma\right)
\end{align}
Then transformations that parameterize all other pure states in the fiber
$\left[  \gamma\right]  _{N}$ in terms of this reference and $\mathfrak{G}%
_{fib}$ non-redundantly can be produced from functions $\phi_{fib}$ defined
as
\begin{align}
\phi_{fib}\left(  \mathsf{z}\right)   &  =\sum_{1\leq\kappa\leq\infty}%
\phi_{fib\kappa}\left(  \left|  \mathsf{z}\right\rangle \left\langle
\mathsf{z}\right|  |\breve{k}_{\kappa}\right)  \left(  D_{refK_{D}}%
^{Nd}|\left|  \mathsf{z}\right\rangle \left\langle \mathsf{z}\right|  \right)
\nonumber\\
\phi_{fib\kappa}  &  =\int k_{\kappa}^{d\ast}\left(  \mathsf{z}\right)
\phi_{fib}\left(  \mathsf{z}\right)  D_{ref}^{N}{}_{K_{D}}\left(
\mathsf{z}\right)  d\mathsf{z};\quad1\leq\kappa\leq\infty\label{refkd}%
\end{align}
where they also satisfy
\begin{align}
\int s_{\iota}^{d\ast}\left(  \mathsf{z}\right)  \phi_{fib}\left(
\mathsf{z}\right)  s_{\kappa}\left(  \mathsf{z}\right)  d\mathsf{z}  &
=0\quad\forall\iota,\kappa\nonumber\\
\int k_{\iota}^{d\ast}\left(  \mathsf{z}\right)  \phi_{fib}\left(
\mathsf{z}\right)  s_{\kappa}\left(  \mathsf{z}\right)  d\mathsf{z}  &
=0\quad\forall\iota,\kappa\nonumber\\
\int s_{\iota}^{d\ast}\left(  \mathsf{z}\right)  \phi_{fib}\left(
\mathsf{z}\right)  k_{\kappa}\left(  \mathsf{z}\right)  d\mathsf{z}  &
=0\quad\forall\iota,\kappa\nonumber\\
\int k_{\iota}^{d\ast}\left(  \mathsf{z}\right)  \phi_{fib}\left(
\mathsf{z}\right)  k_{1}\left(  \mathsf{z}\right)  d\mathsf{z}  &  \neq
0\quad\text{for some }\iota\label{phicond}%
\end{align}
The conditions in eq. $\left(  \ref{phicond}\right)  $ are needed to ensure
that the transformation induced by $\phi_{fib}\left(  \mathsf{z}\right)  $
belong to the defining representation of $\mathcal{G}_{fib}$ on $K_{D}$. We
can the use the diagonal operators $T\left(  \phi_{fib},\omega\right)
=\int\phi_{fib}\left(  \mathsf{z}\right)  e^{i\omega\left(  \mathsf{z}\right)
}\left|  \mathsf{z}\right\rangle \left\langle \mathsf{z}\right|  d\mathsf{z}$,
to define congruence transformations belonging to $\mathfrak{G}_{fib}$ as
\begin{align}
&  \hat{T}_{fib}\left(  \phi_{fib},\omega\right)  Y=T\left(  \phi_{fib}%
,\omega\right)  YT\left(  \phi_{fib},\omega\right)  ^{\dagger}%
=e^{\mathcal{\breve{V}}\left(  \phi_{fib}\right)  +i\breve{\omega}%
}Ye^{\mathcal{\breve{V}}\left(  \phi_{fib}\right)  -i\breve{\omega}%
}\nonumber\\
&  =\left(  1+\breve{\phi}_{fib}\right)  ^{\frac{1}{2}}e^{i\breve{\omega}%
}Y\left(  1+\breve{\phi}_{fib}\right)  ^{\frac{1}{2}}e^{-i\breve{\omega}%
}=\left(  1+\hat{\Phi}_{fib}\right)  \widehat{e^{i\breve{\omega}}}Y
\end{align}
where the integral kernel of $\hat{\Phi}_{fibK_{D}}$ of the projection
$\hat{P}_{K_{D}}\hat{\Phi}_{fib}\hat{P}_{K_{D}}$ of $\hat{\Phi}_{fibK_{D}}$ on
$K_{D}$ is defined as $\Phi_{fibK_{D}}\left(  \mathsf{z}_{1}\mathsf{,z}%
_{2}\right)  =\phi_{fib}\left(  \mathsf{z}_{1}\right)  \delta\left(
\mathsf{z}_{1}-\mathsf{z}_{2}\right)  $. On diagonal operators $Y_{d}%
=Y_{K_{D}}+\Gamma\left(  \gamma\right)  $ the action of these transformations
are
\begin{equation}
\hat{T}_{fib}\left(  \phi_{fib},\omega\right)  Y_{d}=\left(  1+\hat{\Phi
}_{fib}\right)  Y_{d}=Y_{d}+\hat{\Phi}_{fibK_{D}}Y_{K_{D}}=Y_{d}+\breve{\phi
}_{fib}Y_{K_{D}}%
\end{equation}
as $\hat{P}_{S_{D}}\hat{\Phi}_{fib}\hat{P}_{S_{D}}=\hat{P}_{K_{D}}\hat{\Phi
}_{fib}\hat{P}_{S_{D}}=0$. The full intra fiber transformation can thus be
expressed as
\begin{equation}
\hat{T}_{fib}\left(  \phi_{fib},\omega\right)  =e^{\sum_{1\leq\kappa\leq
\infty}\left\{  \Upsilon_{K_{OD}\kappa}\left|  \breve{h}_{\kappa}\right)
\left(  D_{refK_{OD}}^{Nd}\right|  +\Upsilon_{K_{D}\kappa}\left|  \breve
{k}_{\kappa}\right)  \left(  D_{refK_{D}}^{Nd}\right|  \right\}  }
\label{parker}%
\end{equation}
where the coefficients $\Upsilon_{K_{D}\kappa}$ are determined by $\phi_{fib}$
and $\omega$ by the following relationships%
\begin{align}
\mathcal{V}\left(  \phi_{fib}\right)  \left(  \mathsf{z}\right)   &  =\left(
1+\phi_{fib}\left(  \mathsf{z}\right)  \right)  ^{\frac{1}{2}}\nonumber\\
\Upsilon\left(  \mathsf{z}\right)   &  =\mathcal{V}\left(  \phi_{fib}\right)
\left(  \mathsf{z}\right)  +i\omega\left(  \mathsf{z}\right) \nonumber\\
\Upsilon\left(  \mathsf{z}\right)   &  =\sum_{1\leq\kappa\leq\infty}%
\Upsilon_{K_{D}\kappa}\left(  \left|  \mathsf{z}\right\rangle \left\langle
\mathsf{z}\right|  |\breve{k}_{\kappa}\right)  \left(  D_{refK_{D}}%
^{Nd}|\left|  \mathsf{z}\right\rangle \left\langle \mathsf{z}\right|  \right)
\nonumber\\
\Upsilon_{K_{D}\kappa}  &  =\int k_{\kappa}^{d\ast}\left(  \mathsf{z}\right)
\Upsilon\left(  \mathsf{z}\right)  D_{ref}^{N}{}_{K_{D}}\left(  \mathsf{z}%
\right)  d\mathsf{z}%
\end{align}
and the coefficients $\Upsilon_{K_{OD}\kappa}$ are determined by $\phi_{fib}$
and $\omega$. When acting on $D_{ref}^{N}$ the transformation $\hat{T}%
_{fib}\left(  \phi_{fib},\omega\right)  $ produces all other pure states in
$\left[  \gamma\right]  _{N}$ in a unique fashion and thus parameterizes the
pure states in $\left[  \gamma\right]  _{N}$ in a non-redundant manner in
terms of $\left\{  \phi_{fib},\omega\right\}  $. Note that transformations
belonging to $\mathfrak{G}_{fib}$ never change the norm of density operators
as they leave $\Gamma\left(  \gamma\right)  $ unchanged.

The transformations that parameterize $S_{D}$ non-redundantly from a reference
pure state $D_{ref}^{N}$ always may be constructed from selected elements of
$\mathfrak{G}_{den}$ acting on $Y_{ref}$. A complete parameterization of
$S_{D}$ can be achieved in terms of pure states, as given an arbitrary density
$\gamma$ there is \emph{always at least one} pure state that contracts to that
density. If we define $\phi_{den}$ as
\begin{align}
\phi_{den}\left(  \mathsf{z}\right)   &  =\sum_{1\leq\kappa\leq\infty}%
\phi_{den\kappa}\left(  \left|  \mathsf{z}\right\rangle \left\langle
\mathsf{z}\right|  |\breve{s}_{\kappa}\right)  \left(  \Gamma\left(
\gamma\right)  ^{d}|\left|  \mathsf{z}\right\rangle \left\langle
\mathsf{z}\right|  \right) \nonumber\\
\phi_{den\kappa}  &  =\int s_{\kappa}^{d\ast}\left(  \mathsf{z}\right)
\phi_{den}\left(  \mathsf{z}\right)  \Gamma\left(  \gamma\right)  \left(
\mathsf{z}\right)  d\mathsf{z};\quad1\leq\kappa\leq\infty
\end{align}
where they also satisfy
\begin{align}
\int s_{\iota}^{d\ast}\left(  \mathsf{z}\right)  \phi_{den}\left(
\mathsf{z}\right)  s_{1}\left(  \mathsf{z}\right)  d\mathsf{z}  &  \neq
0\quad\text{for some }\iota\nonumber\\
\int k_{\iota}^{d\ast}\left(  \mathsf{z}\right)  \phi_{den}\left(
\mathsf{z}\right)  s_{\kappa}\left(  \mathsf{z}\right)  d\mathsf{z}  &
=0\quad\forall\iota,\kappa\nonumber\\
\int s_{\iota}^{d\ast}\left(  \mathsf{z}\right)  \phi_{den}\left(
\mathsf{z}\right)  k_{\kappa}\left(  \mathsf{z}\right)  d\mathsf{z}  &
=0\quad\forall\iota,\kappa\nonumber\\
\int k_{\iota}^{d\ast}\left(  \mathsf{z}\right)  \phi_{den}\left(
\mathsf{z}\right)  k_{1}\left(  \mathsf{z}\right)  d\mathsf{z}  &
=0\quad\forall\iota,\kappa
\end{align}
We can the use the diagonal operators $T\left(  \phi_{den}\right)  =\int
\phi_{den}\left(  \mathsf{z}\right)  \left|  \mathsf{z}\right\rangle
\left\langle \mathsf{z}\right|  d\mathsf{z}$, to define congruence
transformations%
\begin{align}
\hat{T}_{den}\left(  \phi_{den}\right)  Y  &  =T\left(  \phi_{den}\right)
YT\left(  \phi_{den}\right)  ^{\dagger}=e^{\mathcal{\breve{V}}\left(
\phi_{den}\right)  }Ye^{\mathcal{\breve{V}}\left(  \phi_{den}\right)
}\nonumber\\
&  =\left(  1+\breve{\phi}_{den}\right)  ^{\frac{1}{2}}Y\left(  1+\breve{\phi
}_{den}\right)  ^{\frac{1}{2}}=\left(  1+\hat{\Phi}_{den}\right)  Y
\end{align}
where the integral kernel of $\hat{\Phi}_{denS_{D}}$ of the projection
$\hat{P}_{S_{D}}\hat{\Phi}_{den}\hat{P}_{S_{D}}$ of $\hat{\Phi}_{den}$ on
$S_{D}$ is defined as $\Phi_{den}\left(  \mathsf{z}_{1}\mathsf{,z}_{2}\right)
=\phi_{den}\left(  \mathsf{z}_{1}\right)  \delta\left(  \mathsf{z}%
_{1}-\mathsf{z}_{2}\right)  $. On diagonal operators $Y_{d}=Y_{K_{D}}%
+\Gamma\left(  \gamma\right)  $ this transformation becomes
\begin{equation}
\hat{T}_{den}\left(  \phi_{den}\right)  Y_{d}=\left(  1+\hat{\Phi}%
_{den}\right)  Y_{d}=Y_{d}+\hat{\Phi}_{denS_{D}}\Gamma\left(  \gamma\right)
=Y_{d}+\breve{\phi}_{den}\Gamma\left(  \gamma\right)
\end{equation}
as $\hat{P}_{K_{D}}\hat{\Phi}_{den}\hat{P}_{K_{D}}=\hat{P}_{K_{D}}\hat{\Phi
}_{den}\hat{P}_{S_{D}}=0.$The full intra fiber transformation can thus be
expressed as
\begin{equation}
\hat{T}_{den}\left(  \phi_{den},\omega\right)  =e^{\sum_{1\leq\kappa\leq
\infty}\left\{  \Upsilon_{K_{OD}\kappa}\left|  \breve{h}_{\kappa}\right)
\left(  \breve{Y}_{refK_{OD}}^{d}\right|  +\Upsilon_{S_{D}\kappa}\left|
\breve{s}_{\kappa}\right)  \left(  \Gamma\left(  \gamma\right)  ^{d}\right|
\right\}  } \label{refden}%
\end{equation}
where the coefficients $\Upsilon_{S_{D}\kappa}$ by the following relationships%
\begin{align}
\mathcal{V}\left(  \phi_{den}\right)  \left(  \mathsf{z}\right)   &  =\left(
1+\phi_{den}\left(  \mathsf{z}\right)  \right)  ^{\frac{1}{2}}\nonumber\\
\Upsilon\left(  \mathsf{z}\right)   &  =\mathcal{V}\left(  \phi_{den}\right)
\left(  \mathsf{z}\right) \nonumber\\
\Upsilon\left(  \mathsf{z}\right)   &  =\sum_{1\leq\kappa\leq\infty}%
\Upsilon_{S_{D}\kappa}\left(  \left|  \mathsf{z}\right\rangle \left\langle
\mathsf{z}\right|  |\breve{s}_{\kappa}\right)  \left(  \Gamma\left(
\gamma\right)  ^{d}|\left|  \mathsf{z}\right\rangle \left\langle
\mathsf{z}\right|  \right) \nonumber\\
\Upsilon_{K_{D}\kappa}  &  =\int\breve{s}_{\kappa}^{d\ast}\left(
\mathsf{z}\right)  \Upsilon\left(  \mathsf{z}\right)  \Gamma\left(
\gamma\right)  \left(  \mathsf{z}\right)  d\mathsf{z}%
\end{align}
and the coefficients $\Upsilon_{K_{OD}\kappa}$ are determined by $\phi_{den}$.
When acting on $D_{ref}^{N}$ the transformation $\hat{T}_{den}\left(
\phi_{den}\right)  $ produces a unique set of pure states, one in each
equivalence class. Thus $\hat{T}_{den}\left(  \phi_{den}\right)  $ acting on
the reference pure state $D_{ref}^{N}$ parameterizes $S_{D}$ in a
non-redundant manner in terms of $\left\{  \phi_{den}\right\}  $. It is
possible by left multiplying $\hat{T}_{den}\left(  \phi_{den}\right)  $ by a
transformation $\hat{T}_{fib}\left(  \phi_{fib}\left(  \phi_{den}\right)
,\omega\left(  \phi_{den}\right)  \right)  $ uniquely determined by $\hat
{T}_{den}\left(  \phi_{den}\right)  $, that a path of pure reference states
$D_{ref}^{N}\left(  \gamma\right)  $ of a particular type (i.e. Independent
Particle, Antisymmetrized Geminal Power etc..) and obtain a more specific
non-redundant parameterization of $S_{D}$.

It is further possible to choose a subset of the transformations given in eq.
$\left(  \ref{refden}\right)  $ that produce normalized pure states, (even
from an unnormalized $D_{ref}^{N}$), by letting $\Upsilon_{S_{D}1}$ depend on
$\left\{  \Upsilon_{S_{D}\kappa};2\leq\kappa\leq\infty\right\}  $

\section{Production of Arbitrary Pure States from a Reference Pure State}

\noindent Here we prove that diagonal transformations, in the sense of
Appendix G, can produce any pure state from a reference pure state. The
equivalent mathematical statement is that such transformations are transitive
on the set of pure states.

Associated with any arbitrarily chosen pure state $|\Psi\rangle\,\in
\,\mathcal{H}^{N}\,=\,\mathcal{H}^{N_{e}}\otimes\mathcal{H}^{N_{n}}$ there is
a wavefunction
%\begin{gather}
%\Psi ({\bf x}_{1},\ldots ,{\bf x}_{N_{e}}{\bf X}_{1},\ldots ,{\bf X}%
%_{N_{n}})=\langle {\bf x}_{1},\ldots ,{\bf x}_{N_{e}}{\bf X}_{1},\ldots ,%
%\bf X}_{N_{n}}|\Psi \rangle  \nonumber \\
%|\Psi ({\bf x}_{1},\ldots ,{\bf x}_{N_{e}}{\bf X}_{1},\ldots ,{\bf X}%
%{N_{n}})|e^{i\theta _{\Psi }({\bf x}_{1},\ldots ,{\bf x}_{N_{e}}{\bf X}%
%{1},\ldots ,{\bf X}_{N_{n}})}  \nonumber \\
%equiv e^{{\cal V}_{\Psi }({\bf x}_{1},\ldots ,{\bf x}_{N_{e}}{\bf X}%
%{1},\ldots ,{\bf X}_{N_{n}})}e^{i\theta _{\Psi }({\bf x}_{1},\ldots ,{\bf x}%
%{N_{e}}{\bf X}_{1},\ldots ,{\bf X}_{N_{n}})}  \label{psipolar}
%end{gather}


Consider $T_{d}\in\mathcal{C}(\mathcal{H}^{N})$, a diagonal transformation in
the sense of Appendix G. Its integral kernel can be written as in eq.
(\ref{Tdexplicit}) which leads to eq. (\ref{Tdcoord}). The latter equation in
combination with the last RHS of (\ref{psipolar}) gives
\begin{equation}
\langle\mathsf{x},\mathsf{X}|T_{d}\Psi\rangle=e^{(\mathcal{V}_{\Psi}+
\mathcal{V})}e^{i(\theta_{\Psi}+\omega)}%
\end{equation}
with the obvious string of arguments in compressed notation or surpressed for clarity.

If one chooses an arbitrary $|\chi\rangle\,\in\,\mathcal{H}^{N}$, then the
choice
\begin{align}
\mathcal{V}  &  =\mathcal{V}_{\chi}-\mathcal{V}_{\Psi}\nonumber\\
\omega &  =\theta_{\chi}-\theta_{\Psi} \label{VTthetaT}%
\end{align}
generates that $|\chi\rangle$ from the initial $|\Psi\rangle$, itself
arbitrarily chosen, via the transformation $T_{d}$. Since all physical
densities correspond to pure states \cite{Harriman81}, the result
(\ref{VTthetaT}) means that all densities (fibers) can be reached given a pure
state corresponding to a reference density (fiber) via transformations of the
type developed in Appendix G.

\section{Partitioned Equation of Motion in the $\Upsilon$ Parameterization}

This Appendix tracks the reasoning given in Appendix B but for the complex
analytic parameterization, $\Upsilon=\mathcal{V}+i\omega$ of the fibers, eq.
(\ref{Ffactor}). In this parameterization, the equations of motion that follow
from eq. (\ref{dEdlambda}) - (\ref{Eoft}) can be written in matrix partitioned
form as
\begin{equation}
\left[
\begin{array}
[c]{ccc}%
\eta_{\mathcal{\Upsilon}\,\mathcal{\Upsilon}}(t) & \eta_{\mathcal{\Upsilon
}\,\mathcal{\Upsilon}^{\ast}}(t) & \eta_{\mathcal{\Upsilon}\,\gamma}(t)\\
\eta_{\mathcal{\Upsilon}^{\ast}\,\mathcal{\Upsilon}}(t) & \eta
_{\mathcal{\Upsilon}^{\ast}\,\mathcal{\Upsilon}^{\ast}}(t) & \eta
_{\mathcal{\Upsilon}^{\ast}\,\gamma}(t)\\
\eta_{\gamma\,\mathcal{\Upsilon}}(t) & \eta_{\gamma\,\mathcal{\Upsilon}^{\ast
}}(t) & \eta_{\gamma\,\gamma}(t)
\end{array}
\right]  \left[
\begin{array}
[c]{c}%
\mathcal{\Upsilon}\\
\mathcal{\Upsilon}^{\ast}\\
\dot{\gamma}%
\end{array}
\right]  =\left[
\begin{array}
[c]{c}%
\nabla_{\mathcal{\Upsilon}}\mathcal{E}\left(  t\right) \\
\nabla_{\mathcal{\Upsilon}^{\ast}}\mathcal{E}\left(  t\right) \\
\nabla_{\gamma}\mathcal{E}\left(  t\right)
\end{array}
\right]  \label{3x3partmat}%
\end{equation}
where $^{\ast}$ denotes complex conjugation. These equations may be written
more compactly as
\begin{equation}
\left[
\begin{array}
[c]{cc}%
\eta_{\mathbf{\digamma\digamma}}\left(  t\right)  & \eta_{\mathbf{\digamma
}\gamma}\left(  t\right) \\
\eta_{\gamma\mathbf{\digamma}}\left(  t\right)  & \eta_{\gamma\gamma}\left(
t\right)
\end{array}
\right]  \left[
\begin{array}
[c]{c}%
\mathbf{\dot{\digamma}}\\
\dot{\gamma}%
\end{array}
\right]  =\left[
\begin{array}
[c]{c}%
\nabla_{\mathbf{\digamma}}\mathcal{E}\left(  t\right) \\
\nabla_{\gamma}\mathcal{E}\left(  t\right)
\end{array}
\right]  \label{partmat2}%
\end{equation}
with
\begin{align}
&  \mathbf{\digamma}=\left(
\begin{array}
[c]{c}%
\mathcal{\Upsilon}\\
\mathcal{\Upsilon}^{\ast}%
\end{array}
\right)  \,;\;\eta_{\digamma\digamma}\left(  t\right)  =\left(
\begin{array}
[c]{cc}%
\eta_{\mathcal{\Upsilon\Upsilon}} & \eta_{\mathcal{\Upsilon\Upsilon}^{\ast}}\\
\eta_{\mathcal{\Upsilon}^{\ast}\mathcal{\Upsilon}} & \eta_{\mathcal{\Upsilon
}^{\ast}\mathcal{\Upsilon}^{\ast}}%
\end{array}
\right)  \,;\;\eta_{\mathbf{\digamma}\gamma}\left(  t\right)  =\left(
\begin{array}
[c]{c}%
\eta_{\mathcal{\Upsilon}^{\ast}\gamma}\\
\eta_{\mathcal{\Upsilon}^{\ast}\gamma}%
\end{array}
\right) \nonumber\\
&  \eta_{\gamma\mathbf{\digamma}}\left(  t\right)  =\left(
\begin{array}
[c]{cc}%
\eta_{\gamma\mathcal{\Upsilon}} & \eta_{\gamma\mathcal{\Upsilon}^{\ast}}%
\end{array}
\right)  \,;\;\nabla_{\digamma}\mathcal{E}\left(  t\right)  =\left(
\begin{array}
[c]{c}%
\nabla_{\mathcal{\Upsilon}}\mathcal{E}\left(  t\right) \\
\nabla_{\mathcal{\Upsilon}^{\ast}}\mathcal{E}\left(  t\right)
\end{array}
\right)
\end{align}


In this $\mathbf{{\digamma}}$, $\gamma$ notation the formal structure of the
EOM eq. (\ref{partmat2}) is exactly as in eq. (\ref{partmat1}) of Appendix B
with corresponding definitions for $S$ and $\mathcal{E}$. In parallel with the
development leading to eq. (\ref{cdotgammadot}), it follows that
\begin{align}
\dot{\mathbf{\digamma}}(t)  &  =\eta_{\mathbf{\digamma\,\digamma}}%
^{-1}(t)\left[  \nabla_{\mathbf{\digamma}}\mathcal{E}(t)-\eta
_{\mathbf{\digamma}\,\gamma}(t)\dot{\gamma}\right] \nonumber\\
\dot{\gamma}(t)  &  =\left[  \eta_{\gamma\,\gamma}(t)-\eta_{\gamma
\,\mathbf{\digamma}}(t)\eta_{\mathbf{\digamma}\,\mathbf{\digamma}}^{-1}%
(t)\eta_{\mathbf{\digamma}\,\gamma}(t)\right]  ^{-1}\nonumber\\
&  \times\left[  \nabla_{\gamma}\mathcal{E}(t)-\eta_{\gamma\,\mathbf{\digamma
}}(t)\eta_{\mathbf{\digamma}\,\mathbf{\digamma}}^{-1}(t)\nabla
_{\mathbf{\digamma}}\mathcal{E}(t)\right]
\end{align}
Under the assumption that $\eta_{\mathbf{\digamma\digamma}}^{-1}(t)$ exists,
denote the solution of the first of these for the correct density $\gamma
_{\#}(t)$ by $\mathbf{\digamma}_{\#}(t)$. Then the exact equation of motion
for the temporal development of the density is
\begin{equation}
\dot{\gamma}(t)=\mathcal{N}_{1}(\gamma(t))\left[  \nabla_{\gamma}%
\,\mathcal{E}(\gamma(t))-\mathcal{N}_{2}(\gamma(t))\right]
\end{equation}
The functionals are defined as
%\begin{align}%
\begin{align}
\mathcal{N}_{1}(\gamma(t))  &  =\left[  \eta_{\gamma\,\gamma}(\mathbf{\zeta
}_{\#}(t)\,,\gamma(t))-\eta_{\gamma\,\mathbf{\digamma}}(\mathbf{\zeta}%
_{\#}(t)\,,\gamma(t))\eta_{\mathbf{\digamma}\,\mathbf{\digamma}}%
^{-1}(\mathbf{\zeta}_{\#}(t)\,,\gamma(t))\right.  \times\nonumber\\
&  \qquad\qquad\qquad\qquad\qquad\qquad\qquad\qquad\qquad\qquad\left.
\eta_{\mathbf{\digamma}\,\gamma}(\mathbf{\zeta}_{\#}(t)\,,\gamma(t))\right]
^{-1}\\
\mathcal{N}_{2}\left(  \gamma(t)\right)   &  = \eta_{\gamma\,\mathbf{\digamma
}}(\mathbf{\zeta}_{\#}(t)\,,\gamma(t))\eta_{\mathbf{\digamma}%
\,\mathbf{\digamma}}^{-1}(\mathbf{\zeta}_{\#}(t)\,,\gamma(t))\nabla
_{\mathbf{\digamma}}\mathcal{E}(\gamma(t)).\quad\\
\mathcal{E}(\gamma(t))  &  =\frac{\left\langle \Psi(\mathbf{\zeta}%
_{\#}(t),\gamma(t))|H|\Psi(\mathbf{\zeta}_{\#}(t),\gamma(t))\right\rangle
}{\left\langle \Psi(\mathbf{\zeta}_{\#}(t),\gamma(t))|\Psi(\mathbf{\zeta}%
_{\#}(t),\gamma(t))\right\rangle }%
\end{align}


\begin{thebibliography}{99}                                                                                               %


\bibitem {Teil92}J.~Teilhaber, Phys. Rev. B \textbf{46}, 12990 (1992).

\bibitem {Stener95}M.~Stener, P.~Decleva, and A.~Lisini, J. Phys. B [At. Mol.
Opt. Phys.] \textbf{28}, 4973 (1995).

\bibitem {Saalmann96}U.~Saalmann and R.~Schmidt, Z. Phys. D \textbf{38}, 153 (1996).

\bibitem {Tong97}X.M.~Tong and S-I.~Chu, Phys. Rev. A \textbf{55}, 3406 (1997).

\bibitem {Serra97a}L.~Serra and A.~Rubio, Z. Physik D \textbf{40}, 262 (1997).

\bibitem {Ullrich97}C.A.~Ullrich, P.G.~Reinhard, and E.~Suraud, J. Phys.B [At.
Mol. Optical Phys.] \textbf{30}, 5043 (1997).

\bibitem {AG97}A.~G\"orling, Phys. Rev. B \textbf{55}, 2630 (1997).

\bibitem {Vignale97}G.~Vignale, C.A.~Ullrich, and S.~Conti, Phys. Rev. Lett.
\textbf{79}, 4878 (1997).

\bibitem {Calvayrac98}F.~Calvayrac, A.~Domps, P.G.~Reinhard, E.~Suraud, and
C.A.~Ullrich, Euro. Phys. J. \textbf{4}, 207 (1998).

\bibitem {Reinhard98}P.G.~Reinhard, E.~Suraud, and C.A.~Ullrich, Euro. Phys.
J. \textbf{1}, 303 (1998).

\bibitem {Ullrich98}C.A.~Ullrich, P.G.~Reinhard, and E.~Suraud, Phys. Rev. A
\textbf{57}, 1938 (1998).

\bibitem {Telnov98}D.A.~Telnov and S-I.~Chu, Phys. Rev. A \textbf{58}, 4749 (1998).

\bibitem {Tong98}X.-M.~Tong and S.-I.~Chu, Phys. Rev. A \textbf{57}, 452 (1998).

\bibitem {Campillo99}I.~Campillo, A.~Rubio, and J.M.~Pitarke, Phys. Rev.B
\textbf{59}, 12188 (1999).

\bibitem {Vasiliev99}I.~Vasiliev, S.~Ogut, and J.R.~Chelikowsky, Phys. Rev.
Lett. \textbf{82}, 1919 (1999).

\bibitem {Knospe99}O.~Knospe, J.~Jellinek, U.~Saalmann, and R.~Schmidt, Eur.
Phys. J. D \textbf{5}, 1 (1999) and references therein.

\bibitem {Hessler}P.~Hessler, J.~Park, and K.~Burke, Phys. Rev. Lett.
\textbf{82}, 378 (1999).

\bibitem {Gross96}E.K.U.~Gross, J.F.~Dobson, and M.~Petersilka, in
\textit{Density Functional Theory II - Topics in Current Chemistry 181}, R.F.
Nalewajski ed. (Springer, Berlin and Heidelberg, 1996), 81.

\bibitem {Casida95}M.~Casida in D.P. Chong ed., \textit{Recent Advances in
Density Functional Methods} (World Scientific, Singapore, 1995), 155.

\bibitem {Bartolotti81}L.J.~Bartolotti, Phys. Rev. A \textbf{24}, 1661 (1981).

\bibitem {DebGhosh82}B.M.~Deb and S.K.~Ghosh, J. Chem. Phys. \textbf{77}, 342 (1982).

\bibitem {KD86}H.~Kohl and R.M.~Dreizler, Phys. Rev. Lett. \textbf{56}, 1993 (1986).

\bibitem {RG84}E.~Runge and E.K.U.~Gross, Phys. Rev. Lett. \textbf{52}, 997 (1984).

\bibitem {HK}P.~Hohenberg and W.~Kohn, Phys. Rev. \textbf{136}, B864 (1964).

\bibitem {Xu85}B-X.~Xu and A.K.~Rajagopal, Phys. Rev. A \textbf{31}, 2682 (1985).

\bibitem {Dhara87}A.K.~Dhara and S.K.~Ghosh, Phys. Rev. A \textbf{35}, 442 (1987).

\bibitem {Raj94A}A.K.~Rajagopal, Phys. Lett. A \textbf{193}, 363 (1994).

\bibitem {RajBuot94B}A.K.~Rajagopal and F.A.~Buot, Phys. Rev. E \textbf{50},
721 (1994).

\bibitem {LiTong86}T-C.~Li and P-q. Tong, Phys. Rev. A \textbf{34}, 529 (1986).

\bibitem {Qian98}A.~Qian and V.~Sahni, Phys. Lett. A \textbf{247}, 303 (1998).

\bibitem {Chermyak00}V.~Chernyak and S.~Mukamel J. Chem. Phys. \textbf{112},
3572 (2000).

\bibitem {DeumensA}E.~Deumens, A.~Diz, R.~Longo, and Y.~{\"O}hrn, Rev. Mod.
Phys. \textbf{66}, 917 (1994).

\bibitem {DeumensB}E.~Deumens and Y.~{\"O}hrn, J. Chem. Soc., Faraday Trans.
\textbf{93}, 919 (1997) and refs.~therein.

\bibitem {BLWSBT99a}B.~Weiner and S.B.~Trickey, Adv. Quantum Chemistry
\textbf{35} (Academic, San Diego, 1999), 217.

\bibitem {BLWSBT00a}B.~Weiner and S.B.~Trickey, J. Molec. Struct. (Theochem)
\textbf{501}-\textbf{02}, 65 (2000).

\bibitem {Levy79}M.~Levy, Proc. Natl. Acad. Sci. USA \textbf{76}, 6062 (1979).

\bibitem {LevyPerdew85}M.~Levy and J.P.~Perdew in R.M.~Dreizler and J.~da
Providencia eds., \textit{Density Functional Methods in Physics} (Plenum,
N.Y., 1985), 11.

\bibitem {Lieb}E.H.~Lieb, Internat. J. Quantum. Chem. \textbf{24}, 243 (1983).

\bibitem {PerdewSavin}J.P.~Perdew, A.~Savin, and K.~Burke, Phys. Rev. A
\textbf{51}, 4531 (1995).

\bibitem {BLWSBT98}B.~Weiner and S.B.~Trickey, Internat. J. Quantum Chem.
\textbf{69}, 451 (1998).

\bibitem {SavinSem2}A.~Savin, in J.M. Seminario ed., \textit{Recent
Developments and Applications of Modern Density Functional Theory} (Elsevier,
Amsterdam, 1996), 351.

\bibitem {spinpol}If needed, an approximation corresponding to ordinary
spin-polarized DFT can be extracted from the fully spin-dependent theory by
neglecting the contributions of the off-diagonal elements of the one-electron
space-spin density kernel
\begin{equation}
\left[
\begin{array}
[c]{cc}%
D_{e,\alpha\alpha}^{1}\left(  \mathbf{r,r}^{\prime}\right)  & D_{e,\alpha
\beta}^{1}\left(  \mathbf{r,r}^{\prime}\right) \\
D_{e,\beta\alpha}^{1}\left(  \mathbf{r,r}^{\prime}\right)  & D_{e,\beta\beta
}^{1}\left(  \mathbf{r,r}^{\prime}\right)
\end{array}
\right]  \approx\left[
\begin{array}
[c]{cc}%
D_{e,\alpha\alpha}^{1}\left(  \mathbf{r,r}^{\prime}\right)  & 0\\
0 & D_{e,\alpha\alpha}^{1}\left(  \mathbf{r,r}^{\prime}\right)
\end{array}
\right] \nonumber
\end{equation}
We have not pursued this.

\bibitem {KreibGross}T. Kreibich and E.K.U. Gross, Phys. Rev. Lett
\textbf{86}, 2984 (2001).

\bibitem {Sjoqvist}E.~Sj\"oqvist and M.~Hedstr\"om, Phys. Rev. A \textbf{56},
3417 (1997).

\bibitem {Bruhat}\textit{Analysis, Manifolds, and Physics}, Y.~Choquet-Bruhat
and C. DeWitt-Morette with M. Dillard-Bleick, Revised Edition (North-Holland,
Amsterdam, 1982) 350 ff.

\bibitem {affinebundleref}\textit{Encyclopedic Dictionary of Mathematics},
K.~It\^o ed. (MIT Press, Cambridge, 1996) \textbf{1}, 569.

\bibitem {Yref}For physical densities the existence of arbitrarily many
choices of $Y_{ref}$ is guaranteed by Harriman's theorem \cite{Harriman81}.
For the generalization to trace class operators, the guarantee follows from
the fact that any operator $Y$ always can be decomposed into four positive
operators $p_{j}$ summed as
\[
Y=p_{1}-p_{2}+i(p_{3}-p_{4})
\]
Under the contraction map $\Xi$, each of these yields a positive density (not
necessarily normed to the same value but that is no matter), hence each of the
$p_{j}$ can be constructed by application of Harriman's theorem.

\bibitem {Harriman81}J. Harriman, Phys. Rev. A \textbf{24}, 680 (1981).

\bibitem {diffbundle}Note that here we could consider a subtly different fiber
bundle with base the set of reference elements $\{Y_{ref}\}$ and fiber
elements
\[
Y=\sum_{i=1}^{d}\alpha_{i}\xi_{i}
\]
We use this fiber bundle extensively in section V.

\bibitem {KramSar}P.~Kramer and M.~Saraceno, \textit{Geometry of the Time
Dependent Variational Principle in Quantum Mechanics} (Springer, New York 1981).

\bibitem {etainv}If $\eta^{-1}$ does not exist, classical equations of motion
nonetheless can be obtained in terms of a reduced number of variables that
produce a non-singular metric. For a discussion of proper treatment of the
zeros of $\eta(t,\, \mathbf{y},\, \mathbf{y}^{\prime}) $ in classical
mechanics, see, for example, \textit{Classical Dynamics A Modern Perspective},
E.C.G.~Sudarshan and N.~Mukunda (John Wiley and Sons, NY, 1974) 112 ff.

\bibitem {BLW94}B.~Weiner, E.~Deumens, and Y.~\"{O}hrn, J. Math. Phys.
\textbf{35}, 1139 (1994).

\bibitem {congruence}A congruence transform of an operator A is of the form
$TAT^{\dag}$

\bibitem {RefState}Such a state could be of a simple form (e.g. an Independent
Particle State) that leads to straightforward expressions for expectation
values. Or, it could be a correlated state (e.g., the product of an electronic
Antisymmetrized Geminal Power State \cite{Ortiz81} and a variable width wave
packet for nuclear positions) that provides a better basis for approximations
- though at the expense of more complicated expressions for expectation values.

\bibitem {Ortiz81}J.V.~Ortiz, B.~Weiner, and Y.~\"Ohrn, Internat J. Quantum
Chem. S-\textbf{15}, 113 (1981).

\bibitem {Kummer}H.~Kummer, J. Math. Phys. \textbf{8}, 2063 (1967).

\bibitem {Blaizot86}J-P.~Blaizot and G.~Ripka, \textit{Quantum Theory of
Finite Systems} (MIT Press, Cambridge MA, 1986) Chaps. 17 and 18 and
references therein.

\bibitem {Leeuwen99}R.~van Leeuwen, Phys. Rev. Lett. \textbf{82}, 3863 (1999).
\end{thebibliography}


\end{document}
